% 来源：SGA1.pdf，PDF 页 41--64（第 II 讲；末页为空白）。
\SGAchapter{II}{光滑态射：一般性质与微分性质}{Morphismes lisses : généralités, propriétés différentielles}
\SourcePage{41}
对第 I 讲的引用以 I 标示。提醒读者，本讲的环均为\Term{noetherian-ring}{诺特环}，\Term{prescheme}{预概形}均为\Term{locally-noetherian}{局部诺特}的。

\SGAsection{1}{一般性质}
设 $Y$ 是一个\Term{prescheme}{预概形}，$t_1,\ldots,t_n$ 是不定元，置
\[
Y[t_1,\ldots,t_n]=Y\otimes_{\mathbb Z}\mathbb Z[t_1,\ldots,t_n].\tag{1.1}
\]
于是 $Y[t_1,\ldots,t_n]$ 是 $Y$ 上的一个\Term{affine}{仿射} $Y$-\Term{scheme}{概形}，由\Term{quasi-coherent}{拟凝聚}代数\Term{sheaf}{层} $\mathcal O_Y[t_1,\ldots,t_n]$ 定义。因此，给出这个\Term{prescheme}{预概形}在 $Y$ 上的一个\Term{section}{截面}，等价于给出 $\mathcal O_Y$ 的 $n$ 个\Term{section}{截面}，它们就是相应\Term{homomorphism}{同态}下各 $t_i$ 的像。若 $Y'$ 位于 $Y$ 之上，则
\[
Y[t_1,\ldots,t_n]\times_Y Y'=Y'[t_1,\ldots,t_n].\tag{1.2}
\]
这表明，给出一个从 $Y'$ 到 $Y[t_1,\ldots,t_n]$ 的 $Y$-\Term{morphism}{态射}，等价于给出 $\mathcal O_{Y'}$ 的 $n$ 个\Term{section}{截面}。另一方面，利用 $\mathbb Z$ 上多项式环的相应公式，有
\[
\bigl(Y[t_1,\ldots,t_n]\bigr)[t_{n+1},\ldots,t_m]
=Y[t_1,\ldots,t_m].\tag{1.3}
\]
公式 (1.2) 还表明，$Y[t_1,\ldots,t_n]$ 随 $Y$ \Term{functorial}{函子性}地变化。

$Y[t_1,\ldots,t_n]$ 在 $Y$ 上是\Term{finite-type}{有限型}且\Term{flat}{平坦}的。

\begin{SGAresult}{定义}{1.1}
设 $f:X\to Y$ 是一个\Term{morphism}{态射}，使 $X$ 成为 $Y$-\Term{prescheme}{预概形}。若存在整数 $n\geq0$、$x$ 的一个\Term{open-neighborhood}{开邻域} $U$，以及从 $U$ 到 $Y[t_1,\ldots,t_n]$ 的\Term{etale}{平展} $Y$-\Term{morphism}{态射}，则称 $f$ 在 $x\in X$ 处\Term{smooth}{光滑}\footnote{旧称：$f$ 在 $x$ 处是简单的，或 $x$ 是 $f$ 的一个简单点。这种说法在若干场合容易引起混淆，例如单代数、单群，因而不得不弃用。}，也称 $X$ 在 $x$ 处在 $Y$ 上\Term{smooth}{光滑}。若 $f$ 在 $X$ 的每一点都\Term{smooth}{光滑}，则称 $f$（或 $X$）\Term{smooth}{光滑}。环 $A$ 上的代数 $B$ 称为在 $B$ 的\Term{prime-ideal}{素理想} $\mathfrak p$ 处\Term{smooth}{光滑}，是指 $\operatorname{Spec}(B)$ 在点 $\mathfrak p$ 处在 $\operatorname{Spec}(A)$ 上\Term{smooth}{光滑}。
\SourcePage{42}
若 $B$ 在其每个\Term{prime-ideal}{素理想}处都在 $A$ 上\Term{smooth}{光滑}，则称 $B$ 在 $A$ 上\Term{smooth}{光滑}。最后，\Term{local-ring}{局部环}之间的\Term{local-homomorphism}{局部同态} $A\to B$ 称为\Term{smooth}{光滑}的（或称 $B$ 在 $A$ 上\Term{smooth}{光滑}）\footnote{这里最好像 EGA IV 18.6.1 那样，称 $B$ 在 $A$ 上“本质\Term{smooth}{光滑}”。}，是指 $B$ 是某个在 $A$ 上\Term{smooth}{光滑}的\Term{finite-type}{有限型} $A$-代数 $B_1$ 的\Term{localization}{局部化}。
\end{SGAresult}

注意，$X$ 在 $Y$ 上\Term{smooth}{光滑}这一概念在 $X$ 和 $Y$ 上都是局部的；若 $X$ 在 $Y$ 上\Term{smooth}{光滑}，则它在 $Y$ 上是\Term{locally-finite-type}{局部有限型}的。

\begin{SGAresult}{命题}{1.1}
$X$ 中使 $f$ \Term{smooth}{光滑}的点 $x$ 所成的集合是开集。
\end{SGAresult}
这从定义立即可见。

\begin{SGAresult}{推论}{1.2}
若 $B$ 在 $\mathfrak p$ 处在 $A$ 上\Term{smooth}{光滑}，则对 $B$ 中每个包含于 $\mathfrak p$ 的\Term{prime-ideal}{素理想} $\mathfrak q$，$B$ 在 $\mathfrak q$ 处也在 $A$ 上\Term{smooth}{光滑}。
\end{SGAresult}

1.1 还表明，定义 1.1 中最后两个定义在共同适用的范围内一致。

\begin{SGAresult}{命题}{1.3}
\begin{enumerate}
\item[(i)] \Term{etale}{平展}\Term{morphism}{态射}是\Term{smooth}{光滑}的；特别地，\Term{open-immersion}{开浸入}和恒等\Term{morphism}{态射}都是\Term{smooth}{光滑}的。
\item[(ii)] \Term{smooth}{光滑}\Term{morphism}{态射}经\Term{base-change}{基变换}后仍\Term{smooth}{光滑}。
\item[(iii)] 两个\Term{smooth}{光滑}\Term{morphism}{态射}的复合仍\Term{smooth}{光滑}。
\end{enumerate}
\end{SGAresult}
(i) 从定义立即可见；更确切地，有：

\begin{SGAresult}{推论}{1.4}
\Term{etale}{平展} $=$ \Term{quasi-finite}{拟有限} $+$ \Term{smooth}{光滑}。
\end{SGAresult}

(ii) 由\Term{etale}{平展}\Term{morphism}{态射}的相应性质（I 4.6）和投影 $Y[t_1,\ldots,t_n]\to Y$ 的相应性质（见 (1.2)）立即得出。对于 (iii)，\Term{etale}{平展}\Term{morphism}{态射}和形如 $Y[t_1,\ldots,t_n]\to Y$ 的投影分别满足复合封闭性（见 I 4.6 与 (1.3)），再结合证明 (ii) 时用到的两个事实，便形式地得到结论。具体说，设 $Y$ 在 $Z$ 上\Term{smooth}{光滑}，$X$ 在 $Y$ 上\Term{smooth}{光滑}，要证 $X$ 在 $Z$ 上\Term{smooth}{光滑}。可设 $Y$ 在 $Z[t_1,\ldots,t_n]$ 上\Term{etale}{平展}，而 $X$ 在 $Y[s_1,\ldots,s_m]$ 上\Term{etale}{平展}。第一个假设说明 $Y[s_1,\ldots,s_m]$ 在
\[
Z[t_1,\ldots,t_n][s_1,\ldots,s_m]=Z[t_1,\ldots,s_m]
\]
上\Term{etale}{平展}，因而 $X$ 在 $Z[t_1,\ldots,s_m]$ 上\Term{etale}{平展}，证毕。

\textbf{注记 1.5。}定义 1.1 中的整数 $n$ 是唯一确定的，因为立即可见，它就是点 $x$ 在其\Term{fiber}{纤维} $f^{-1}(f(x))$ 中的\Term{local-ring}{局部环}的维数。称它为 $X$ 在 $Y$ 上的“\Term{relative-dimension}{相对维数}”。它对\Term{morphism}{态射}的复合满足可加性。

\SourcePage{43}
\SGAsection{2}{态射光滑性的若干判别准则}
\begin{SGAresult}{定理}{2.1}
设 $f:X\to Y$ 是\Term{locally-finite-type}{局部有限型}\Term{morphism}{态射}，$x\in X$，$y=f(x)$。$f$ 在 $x$ 处\Term{smooth}{光滑}的充要条件是：\textup{(a)} $f$ 在 $x$ 处\Term{flat}{平坦}；\textup{(b)} $f^{-1}(y)$ 在 $x$ 处在 $k(y)$ 上\Term{smooth}{光滑}。
\end{SGAresult}
\begin{proof}
两个\Term{flat}{平坦}\Term{morphism}{态射}的复合是\Term{flat}{平坦}的，而 $Y[t_1,\ldots,t_n]\to Y$ \Term{flat}{平坦}，可知\Term{smooth}{光滑}蕴含\Term{flat}{平坦}；结合 1.3 (ii)，这就证明了必要性。设 (a)、(b) 成立，取 $y$ 的\Term{affine}{仿射}邻域 $V$，其环为 $A$；再取 $x$ 的、位于 $V$ 之上的\Term{affine}{仿射}邻域 $U$，其环为 $B$。把 $U$ 取得足够小，由 (b) 可设存在\Term{etale}{平展}的 $k(y)$-\Term{morphism}{态射}
\[
g:U|_{f^{-1}(y)}\longrightarrow\operatorname{Spec}k[t_1,\ldots,t_n],
\qquad k=k(y),
\]
由 $U|_{f^{-1}(y)}$ 的结构\Term{sheaf}{层}的 $n$ 个\Term{section}{截面} $g_i$ 定义。不难看出，可以设 $g_i$ 来自 $U$ 的结构\Term{sheaf}{层}的\Term{section}{截面}：它们先验只是 $B\otimes_A k=BS^{-1}$ 中的元素，这里 $S=A-\mathfrak p$，$\mathfrak p$ 是 $A$ 中对应于 $y$ 的\Term{prime-ideal}{素理想}；必要时，把各 $g_i$ 同乘以 $k$ 中一个非零元素即可。因此 $g$ 由一个仍记作 $g$ 的\Term{morphism}{态射}
\[
g:U\longrightarrow Y[t_1,\ldots,t_n]
\]
诱导。由 (a)，$U$ 在 $Y$ 上\Term{flat}{平坦}；$Y[t_1,\ldots,t_n]$ 也在 $Y$ 上\Term{flat}{平坦}。又 $g$ 在 $y$ 上方的\Term{fiber}{纤维}之间诱导\Term{etale}{平展}\Term{morphism}{态射}，故由 I 5.8，$g$ 在 $x$ 处\Term{etale}{平展}，证毕。
\end{proof}

\begin{SGAresult}{推论}{2.2}
设 $S$ 是\Term{prescheme}{预概形}，$f:X\to Y$ 是\Term{finite-type}{有限型} $S$-\Term{morphism}{态射}，$Y$ 在 $S$ 上\Term{finite-type}{有限型}且\Term{flat}{平坦}，$x\in X$，$s$ 是 $x$ 在 $S$ 上的投影。$f$ 在 $x$ 处\Term{smooth}{光滑}的充要条件是：$X$ 在 $x$ 处在 $S$ 上\Term{flat}{平坦}（也可要求\Term{smooth}{光滑}），且\Term{fiber}{纤维}间的诱导\Term{morphism}{态射} $f_s:X_s\to Y_s$ 在 $x$ 处\Term{smooth}{光滑}。
\end{SGAresult}
只须证明充分性，它由准则 2.1 和\Term{flatness}{平坦性}准则 I 5.9 得出。

为叙述下一个结果，先“回顾”如下定义。\Term{locally-finite-type}{局部有限型}\Term{morphism}{态射} $f:X\to Y$ 称为在点 $x\in X$ 处\Term{equidimensional}{等维}，是指置 $y=f(x)$ 后，可取 $x$ 的一个\Term{open-neighborhood}{开邻域} $U$，其每个分支都\Term{dominant}{支配} $Y$ 的某个分支，并且对任意 $y'\in Y$，$f^{-1}(y')\cap U$ 的所有\Term{irreducible-component}{不可约分支}都有同一个维数，该维数与 $y'$ 无关。实际上，只对经过 $y$ 的 $Y$ 的各分支的\Term{generic-point}{泛点}及点 $y$ 本身考察这一条件就够了。例如，若 $X$、$Y$ 都整且 $f$ \Term{dominant}{支配}，则该条件是说，$f^{-1}(y)$ 中经过 $x$ 的分支具有“正确的”维数，即\Term{generic-fiber}{泛纤维}的维数；提醒读者，它们的维数总不小于\Term{generic-fiber}{泛纤维}的维数。若 $f$ 在 $x$ 处\Term{equidimensional}{等维}，其\Term{fiber}{纤维}在 $x$ 处的维数为 $n$，且
\[
g:U\longrightarrow Y'=Y[t_1,\ldots,t_n]
\]
是 $x$ 的某个邻域上的 $Y$-\Term{morphism}{态射}，在 $y$ 上方的\Term{fiber}{纤维}之间诱导的\Term{morphism}{态射}在 $x$ 处\Term{quasi-finite}{拟有限}（等价地，$g$ 在 $x$ 处\Term{quasi-finite}{拟有限}），则可以证明，
\SourcePage{44}
$U$ 中每个经过 $x$ 的分支都\Term{dominant}{支配} $Y'$ 的某个\Term{irreducible-component}{不可约分支}。此外，由“\Term{normalization}{正规化}引理”，这样的 $g$ 总存在。反过来，若 $x$ 的某个\Term{open-neighborhood}{开邻域} $U$ 到形如 $Y'=Y[t_1,\ldots,t_n]$ 的 $Y$-\Term{scheme}{概形}存在\Term{quasi-finite}{拟有限}的 $Y$-\Term{morphism}{态射} $g$，并且 $U$ 的每个分支都\Term{dominant}{支配} $Y'$ 的某个分支，则 $f$ 在 $x$ 处\Term{equidimensional}{等维}。于是有：

\begin{SGAresult}{命题}{2.3}
设 $f:X\to Y$ 是\Term{locally-finite-type}{局部有限型}\Term{morphism}{态射}，$x\in X$，$y=f(x)$，并设 $\mathcal O_y$ \Term{normal}{正规}。则 $f$ 在 $x$ 处\Term{smooth}{光滑}的充要条件是：$f$ 在 $x$ 处\Term{equidimensional}{等维}，且 $f^{-1}(y)$ 在 $x$ 处在 $k(y)$ 上\Term{smooth}{光滑}。
\end{SGAresult}
\begin{proof}
从定义立即可见，\Term{smooth}{光滑}\Term{morphism}{态射}是\Term{equidimensional}{等维}的。注意，\Term{finite-type}{有限型}\Term{flat}{平坦}\Term{morphism}{态射}未必在 $x$ 处\Term{equidimensional}{等维}，即使其\Term{fiber}{纤维}在 $x$ 处\Term{irreducible}{不可约}，也未必如此。证明逆命题。因为 $f^{-1}(y)$ 在 $x$ 处在 $k(y)$ 上\Term{smooth}{光滑}，必要时以合适邻域替换 $X$，可设存在 $Y$-\Term{morphism}{态射}
\[
g:X\longrightarrow Y[t_1,\ldots,t_n]=Y',
\]
在 $y$ 上方的\Term{fiber}{纤维}之间诱导\Term{etale}{平展}\Term{morphism}{态射}，因而在 $x$ 处\Term{quasi-finite}{拟有限}。所以 $g$ \Term{unramified}{非分歧}；又因 $f$ 在 $x$ 处\Term{equidimensional}{等维}，$X$ 中经过 $x$ 的各分支都\Term{dominant}{支配} $Y'$ 的某个分支。因此 $g$ 诱导的\Term{homomorphism}{同态} $\mathcal O_{y'}\to\mathcal O_x$ 是\Term{injectivity}{单射}，其中 $y'=g(x)$。这个\Term{homomorphism}{同态}还\Term{unramified}{非分歧}，而 $\mathcal O_{y'}$ 是\Term{normal}{正规}环 $\mathcal O_y[t_1,\ldots,t_n]$ 的\Term{localization}{局部化}，故\Term{normal}{正规}。所以该\Term{homomorphism}{同态}\Term{etale}{平展}（I 9.5 (ii)）。
\end{proof}

\textbf{注记 2.4。}上述结论中，“$\mathcal O_y$ \Term{normal}{正规}”可以换成更弱的条件：“$Y$ 在 $y$ 处\Term{geometrically-unibranch}{几何单分支}”（见 I 11），因为 I 9.5 在这个假设下仍然成立。借此机会也指出：若整\Term{local-ring}{局部环} $A$ 的\Term{residue-field}{剩余域}\Term{algebraically-closed}{代数闭}，则“\Term{analytically-integral}{解析整}”，即 $\widehat A$ 整，蕴含\Term{geometrically-unibranch}{几何单分支}。在任意“良好环”的\Term{category}{范畴}中，逆命题也成立；确切地说，该\Term{category}{范畴}须在通常运算下稳定，而且其中\Term{normal}{正规}\Term{local-ring}{局部环}的\Term{completion}{完备化}仍\Term{normal}{正规}。由 Zariski 的“解析\Term{normal}{正规}性定理”，\Term{affine}{仿射}代数及其\Term{localization}{局部化}组成的\Term{category}{范畴}满足这个条件。\footnote{见 EGA IV 7.8。}

最后，在当前的讨论范围内，再“回顾”一个 Hironaka 的结果\footnote{见 EGA IV 5.12.10。}。它有时可用来判定 $f^{-1}(y)$ 是\Term{reduced}{约化}\Term{scheme}{概形}；换言之，它与过去许多代数几何学家不甚恰当地称为 $f$ 在 $x$ 上方的“无\Term{multiplicity}{重数}\Term{fiber}{纤维}”的对象，即 $f^{-1}(y)_{\mathrm{red}}$，确实相同：

\SourcePage{45}
\begin{SGAresult}{命题}{2.5}
设 $f:X\to Y$ 是\Term{reduced}{约化}\Term{prescheme}{预概形}之间的\Term{finite-type}{有限型}\Term{dominant}{支配}\Term{morphism}{态射}，$y\in Y$，且 $\mathcal O_y$ \Term{regular}{正则}。设 $f^{-1}(y)$ 的所有分支的\Term{multiplicity}{重数}都为 $1$（定义见下文），且 $f^{-1}(y)_{\mathrm{red}}$ \Term{normal}{正规}。则 $f^{-1}(y)$ \Term{reduced}{约化}，从而\Term{normal}{正规}；$X$ 在 $f^{-1}(y)$ 的每一点都\Term{normal}{正规}；并且 $X$ 在 $f^{-1}(y)$ 的每一点都在 $Y$ 上\Term{flat}{平坦}。
\end{SGAresult}

$f^{-1}(y)$ 的一个分支 $Z$ 称为\Term{multiplicity}{重数}为 $1$，是指设 $x$ 为 $Z$ 的\Term{generic-point}{泛点}后：\textup{(i)} $\dim\mathcal O_x=\dim\mathcal O_y$，即 $Z$ 不是“多余分支”，也就是说，它的维数并非“过大”；\textup{(ii)} $\mathcal O_x$ 的\Term{maximal-ideal}{极大理想}由 $\mathcal O_y$ 的\Term{maximal-ideal}{极大理想}生成。按 $x$ 的选法，后者先验在 $\mathcal O_x$ 中生成一个\Term{ideal-of-definition}{定义理想}。

结合 2.3 或 2.1，得到：
\begin{SGAresult}{推论}{2.6}
设 $f:X\to Y$ 是\Term{reduced}{约化}\Term{prescheme}{预概形}之间的\Term{finite-type}{有限型}\Term{dominant}{支配}\Term{morphism}{态射}，$y\in Y$，$\mathcal O_y$ \Term{regular}{正则}。则 $f$ 在 $X$ 中位于 $y$ 上方的各点\Term{smooth}{光滑}的充要条件是：$f^{-1}(y)$ 的各分支的\Term{multiplicity}{重数}均为 $1$，且 $f^{-1}(y)_{\mathrm{red}}$ 在 $k(y)$ 上\Term{smooth}{光滑}。
\end{SGAresult}

过去，人们主要在 $Y$ 是某个\Term{discrete-valuation-ring}{离散赋值环} $A$ 的谱时考虑这种情况，并常说成“若 $X$ 关于所给赋值的\Term{reduced}{约化}是良好的”之类。那时，$X$ 还表示 $\mathbb P^n_K$ 的一个闭子\Term{scheme}{概形}——如果当时的说法还配得上“子\Term{scheme}{概形}”这个名字的话——这里 $K$ 是 $A$ 的\Term{field-of-fractions}{分式域}。由于缺少合适的语言，人们几乎未能看出，一个对象“定义在 $A$ 上”（而非仅仅定义在 $K$ 上）具有更内在的意义。

\SGAsection{3}{保持性质}
\begin{SGAresult}{命题}{3.1}
设 $f:X\to Y$ 是\Term{morphism}{态射}，$x\in X$，$y=f(x)$。假设 $f$ 在 $x$ 处\Term{smooth}{光滑}。则 $\mathcal O_x$ \Term{reduced}{约化}（分别地，\Term{regular}{正则}、\Term{normal}{正规}）的充要条件是 $\mathcal O_y$ 具有相应性质。
\end{SGAresult}

当 $X$ 形如 $Y[t_1,\ldots,t_n]$ 时，这一结论已知；对\Term{etale}{平展}\Term{morphism}{态射}，它已在第 I 讲第 9 节证明。一般情形便由定义 1.1 立即得到。

其余保持性质来自\Term{flatness}{平坦性}本身，或来自以下事实：$X$ 在某个形如 $Y[t_1,\ldots,t_n]$ 的 $Y$-\Term{prescheme}{预概形}上局部\Term{quasi-finite}{拟有限}且\Term{flat}{平坦}；我们也将称这种情况为 $X$ 在 $Y$ 上是\Term{cohen-macaulay}{Cohen--Macaulay}的。这里不再详述，只指出最后这个事实给出
\[
\dim\mathcal O_x=\dim\mathcal O_y+n-d,\qquad
\operatorname{prof}\mathcal O_x=\operatorname{prof}\mathcal O_y+n-d,\tag{3.1}
\]
\SourcePage{46}
其中 $n$ 是 $f$ 的\Term{fiber}{纤维}在 $x$ 处的维数，$d$ 是 $k(x)$ 在 $k(y)$ 上的\Term{transcendence-degree}{超越次数}。因此，置 $\operatorname{coprof}=\dim-\operatorname{prof}$，有\footnote{这些公式见 EGA IV 6.1 与 6.3。}
\[
\operatorname{coprof}\mathcal O_x=\operatorname{coprof}\mathcal O_y.\tag{3.2}
\]
例如，$\mathcal O_x$ 是\Term{cohen-macaulay}{Cohen--Macaulay}的（分别地，无\Term{embedded-component}{嵌入分支}）当且仅当 $\mathcal O_y$ 亦然。

\SGAsection{4}{光滑态射的微分性质}
为简化讨论，我们主要限于一阶微分计算；对高阶情形只作简要说明，其结果同样简单。

$Y$-\Term{prescheme}{预概形} $X$ 的一阶\Term{differential-sheaf}{微分层} $\Omega^1_{X/Y}$ 的定义见第 I 讲第 1 节。设 $X$、$Y$ 是 $S$-\Term{prescheme}{预概形}，结构\Term{morphism}{态射} $f:X\to Y$ 是 $S$-\Term{morphism}{态射}。则 $f$ 定义一个与 $f$ 相容的模\Term{homomorphism}{同态}
\[
f^*:\Omega^1_{Y/S}\longrightarrow\Omega^1_{X/S}.\tag{4.1}
\]
换言之，$\Omega^1_{X/S}$ 对 $S$-\Term{prescheme}{预概形} $X$ 是反变的。此外，(4.1) 等价于 $X$ 上的一个模\Term{homomorphism}{同态}
\[
f^*(\Omega^1_{Y/S})\longrightarrow\Omega^1_{X/S},\tag{4.1bis}
\]
没有更好的记号，仍以 $f^*$ 表示。它嵌入如下典范\Term{exact-sequence}{正合列}：
\[
f^*(\Omega^1_{Y/S})\longrightarrow\Omega^1_{X/S}
\longrightarrow\Omega^1_{X/Y}\longrightarrow0.\tag{4.2}
\]
这些\Term{homomorphism}{同态}均由两个要求确定：具有局部性，从而可归结为\Term{affine}{仿射}情形；与算子 $d$ 可交换。(4.2) 的\Term{exactness}{正合性}是经典而直接的；在\Term{affine}{仿射}情形，它写成对应于 $A$-代数\Term{homomorphism}{同态} $B\to C$ 的下述\Term{exact-sequence}{正合列}：
\[
\Omega^1_{B/A}\otimes_B C\longrightarrow\Omega^1_{C/A}
\longrightarrow\Omega^1_{C/B}\longrightarrow0.\tag{4.2bis}
\]

\begin{SGAresult}{引理}{4.1}
设 $f:X\to Y$ 是 $S$-\Term{prescheme}{预概形}的\Term{morphism}{态射}。若 $f$ \Term{unramified}{非分歧}（分别地，\Term{etale}{平展}），则 $f^*(\Omega^1_{Y/S})\to\Omega^1_{X/S}$ 是\Term{surjectivity}{满射}（分别地，\Term{isomorphism}{同构}）。在“\Term{unramified}{非分歧}”的情形，若假设 $f$ \Term{locally-finite-type}{局部有限型}，则逆命题也成立。
\end{SGAresult}

\SourcePage{47}
\Term{unramified}{非分歧}的情形由\Term{exact-sequence}{正合列} (4.2) 和 I 3.1 得出；它也和\Term{etale}{平展}情形一样，可直接看出。考虑\Term{commutative-diagram}{交换图}
\[
\begin{tikzcd}
X \arrow[r,"\Delta_{X/Y}"] & X\times_Y X \arrow[r]\arrow[d]
  & X\times_S X \arrow[d]\\
 &Y \arrow[r,"\Delta_{Y/S}"] &Y\times_S Y.
\end{tikzcd}
\]
这里 $X\times_Y X$ 等同于 $Y$ 与 $X\times_S X$ 在 $Y\times_S Y$ 上的\Term{fiber-product}{纤维积}。由于 $f$ \Term{unramified}{非分歧}，$X\to X\times_Y X$ 是\Term{open-immersion}{开浸入}，故它与 $X\times_Y X\to X\times_S X$ 复合而成的\Term{immersion}{浸入} $\Delta_{X/S}$ 的\Term{conormal-sheaf}{余法层}，\Term{isomorphism}{同构}于\Term{immersion}{浸入} $X\times_Y X\to X\times_S X$ 的\Term{conormal-sheaf}{余法层}在 $X$ 上的\Term{inverse-image}{逆像}。另一方面，$X\to Y$ \Term{etale}{平展}，因而\Term{flat}{平坦}，故 $X\times_S X\to Y\times_S Y$ \Term{flat}{平坦}。所以\Term{immersion}{浸入} $X\times_Y X\to X\times_S X$ 的\Term{conormal-sheaf}{余法层}，\Term{isomorphism}{同构}于\Term{immersion}{浸入} $Y\to Y\times_S Y$ 的\Term{conormal-sheaf}{余法层}的\Term{inverse-image}{逆像}，也就是 $\Omega^1_{Y/S}$ 的\Term{inverse-image}{逆像}。结论随之得到。

\begin{SGAresult}{引理}{4.2}
设 $X=Y[t_1,\ldots,t_n]$，其中 $Y$ 是 $S$-\Term{prescheme}{预概形}。则典范\Term{homomorphism}{同态}列
\[
0\longrightarrow f^*(\Omega^1_{Y/S})\longrightarrow\Omega^1_{X/S}
\longrightarrow\Omega^1_{X/Y}\longrightarrow0
\]
\Term{exact}{正合}，且 $\Omega^1_{X/Y}$ 是以各 $d_{X/Y}t_i$ 为基的\Term{free-module}{自由模}。
\end{SGAresult}
核对完全在\Term{affine}{仿射}情形下进行，并且十分直接。注意，(4.2) 的\Term{exactness}{正合性}已经知道。

综合以上两个结论和定义 1.1，得到：
\begin{SGAresult}{定理}{4.3}
设 $f:X\to Y$ 是 $S$-\Term{prescheme}{预概形}之间的\Term{smooth}{光滑}\Term{morphism}{态射}。则：
\begin{enumerate}
\item[(i)] 典范\Term{homomorphism}{同态}列
\[
0\longrightarrow f^*(\Omega^1_{Y/S})\longrightarrow\Omega^1_{X/S}
\longrightarrow\Omega^1_{X/Y}\longrightarrow0
\]
\Term{exact}{正合}。
\item[(ii)] $\Omega^1_{X/Y}$ \Term{locally-free}{局部自由}，它在 $x$ 处的\Term{rank}{秩} $n$ 等于 $f$ 在 $x$ 处的\Term{relative-dimension}{相对维数}。
\end{enumerate}
\end{SGAresult}

\begin{SGAresult}{推论}{4.4}
\Term{homomorphism}{同态} $f^*(\Omega^1_{Y/S})\to\Omega^1_{X/S}$ 是\Term{injectivity}{单射}，其像在 $\Omega^1_{X/S}$ 中局部为\Term{direct-summand}{直和项}。
\end{SGAresult}

设 $u:F\to G$ 是\Term{prescheme}{预概形} $X$ 上的模\Term{homomorphism}{同态}。称它在 $x\in X$ 处\Term{universally-injective}{普遍单射}，是指 $\mathcal O_x$-模\Term{homomorphism}{同态} $F_x\to G_x$ 为\Term{injectivity}{单射}，并且与任意 $\mathcal O_x$-代数作\Term{tensor-product}{张量积}后仍为\Term{injectivity}{单射}；等价地，也可与任意 $\mathcal O_x$-模作\Term{tensor-product}{张量积}。例如，若 $x$ 有一个\Term{open-neighborhood}{开邻域} $U$，使 $u$ 把 $F|_U$ \Term{isomorphism}{同构}地映到 $G|_U$ 的一个\Term{direct-summand}{直和项}，则这已足够。当 $F$、$G$ 在 $x$ 的某个邻域上自由且\Term{finite-type}{有限型}时，该条件也必要。更确切地，此时下列条件等价：
\SourcePage{48}
\begin{enumerate}
\item[(i)] $u$ 在 $x$ 处\Term{injectivity}{单射}，且 $\operatorname{Coker}u$ 在 $x$ 处自由；
\item[(ii)] $x$ 有\Term{open-neighborhood}{开邻域} $U$，使 $u$ 把 $F|_U$ \Term{isomorphism}{同构}地映到 $G|_U$ 的一个\Term{direct-summand}{直和项}；
\item[(iii)] $u$ 在 $x$ 处\Term{universally-injective}{普遍单射}；
\item[(iv)] \Term{homomorphism}{同态} $F_x\otimes k(x)\to G_x\otimes k(x)$，即它在点 $x$ 的\Term{fiber}{纤维}之间诱导的\Term{homomorphism}{同态}，是\Term{injectivity}{单射}；
\item[(v)] 转置\Term{homomorphism}{同态} $\check G\to\check F$ 在 $x$ 处\Term{surjectivity}{满射}；等价地，它在 $x$ 的某个邻域上\Term{surjectivity}{满射}。
\end{enumerate}
这些条件可循环证明。(iv)$\Rightarrow$(v) 由 Nakayama 引理得出；(v)$\Rightarrow$(i)，因为\Term{locally-free}{局部自由}的商\Term{sheaf}{层}必为\Term{direct-summand}{直和项}。几何上，这意味着 $u$ 对应于一个\Term{isomorphism}{同构}，把以 $F$ 为\Term{section}{截面}\Term{sheaf}{层}的\Term{vector-bundle}{向量丛}，\Term{isomorphism}{同构}地映到以 $G$ 为\Term{section}{截面}\Term{sheaf}{层}的\Term{vector-bundle}{向量丛}的一个\Term{subbundle}{子丛}上。当然，仅有 $F\to G$ 为\Term{injectivity}{单射}还不够。

\begin{SGAresult}{推论}{4.5}
设 $f:X\to Y$ 是 $S$-\Term{prescheme}{预概形}之间的\Term{locally-finite-type}{局部有限型}\Term{morphism}{态射}，$x\in X$，$y=f(x)$，$s$ 是 $x$、$y$ 在 $S$ 上的投影。设 $Y$ 在 $y$ 处在 $S$ 上\Term{smooth}{光滑}。下列条件等价：
\begin{enumerate}
\item[(i)] $f$ 在 $x$ 处\Term{smooth}{光滑}；
\item[(ii)] $X$ 在 $x$ 处在 $S$ 上\Term{smooth}{光滑}，且 $f^*(\Omega^1_{Y/S})\to\Omega^1_{X/S}$ 在 $x$ 处\Term{universally-injective}{普遍单射}；也就是说，它在 $x$ 处\Term{injectivity}{单射}，其\Term{cokernel}{余核} $\Omega^1_{X/Y}$ 在 $x$ 处自由。
\end{enumerate}
\end{SGAresult}
必要性来自 1.3 (iii) 和 4.3 (i)、(ii)。证明充分性。因各 $dg$（$g\in\mathcal O_x$）在 $x$ 处生成模 $\Omega^1_{X/Y}$，可取 $g_i$（$1\leq i\leq n$），使 $dg_i$ 在 $(\Omega^1_{X/Y})_x$ 中的像构成该模的一组基。把 $X$ 取得足够小，可设 $g_i$ 来自 $\mathcal O_X$ 的\Term{section}{截面}，从而定义 $Y$-\Term{morphism}{态射}
\[
g:X\longrightarrow Y'=Y[t_1,\ldots,t_n].
\]
利用假设和引理 4.2，不难看出相应\Term{homomorphism}{同态} $g^*(\Omega^1_{Y'/S})\to\Omega^1_{X/S}$ 在 $x$ 处为\Term{bijection}{双射}。因此归结为证明下一个结论。

\begin{SGAresult}{推论}{4.6}
设 $f:X\to Y$ 是\Term{smooth}{光滑} $S$-\Term{prescheme}{预概形}之间的\Term{morphism}{态射}。$f$ 在 $x\in X$ 处\Term{etale}{平展}的充要条件是：\Term{homomorphism}{同态} $f^*(\Omega^1_{Y/S})\to\Omega^1_{X/S}$ 在 $x$ 处为\Term{isomorphism}{同构}。
\end{SGAresult}
必要性由 4.1 已知。由同一引理，这个条件蕴含 $f$ 在 $x$ 处\Term{unramified}{非分歧}。由 2.2，可归结为 $S=\operatorname{Spec}(k)$ 的情形。因为 $Y$ 在 $k$ 上\Term{smooth}{光滑}，故\Term{regular}{正则}，从而\Term{normal}{正规}；由 I 9.5 (ii)，只须证明 $\mathcal O_y\to\mathcal O_x$ \Term{injectivity}{单射}，或者说 $\mathcal O_y$ 与 $\mathcal O_x$ 的维数相同。这两个维数分别是 $\Omega^1_{Y/k}$ 在 $y$ 处和 $\Omega^1_{X/k}$ 在 $x$ 处的\Term{rank}{秩}，按假设二者相同。

\textbf{注记 4.7。}设 $X$、$Y$ 在 $S$ 上\Term{smooth}{光滑}。$f:X\to Y$ 的\Term{smoothness}{光滑性}准则 4.5 (ii) 还可以表述为：在 $X$ 的每一点 $x$，$f$ 相对于基底 $S$ 的切映射是\Term{surjectivity}{满射}。它就是 $f^*(\Omega^1_{Y/S})$ 与 $\Omega^1_{X/S}$ 在 $x$ 的\Term{fiber}{纤维}——两个\Term{finite}{有限}维 $k(x)$-向量空间——之间的\Term{homomorphism}{同态}的转置。
\SourcePage{49}
这种假设十分常见，特别是研究\Term{analytic-space}{解析空间}的人对它很熟悉。他们通常还加上\Term{nonsingular}{非奇异}性假设，即 $X$、$Y$ “在 $\mathbb C$ 上\Term{smooth}{光滑}”，见第 5 节。看来，这只是因为不少几何学家至今仍对代数簇或\Term{analytic-space}{解析空间}的奇点心存畏惧。

下面指出 4.6 的一个特例。
\begin{SGAresult}{推论}{4.8}
设 $X$ 是 $S$-\Term{prescheme}{预概形}，$g:X\to S[t_1,\ldots,t_n]$ 是由 $\mathcal O_X$ 的\Term{section}{截面} $g_i$（$1\leq i\leq n$）定义的 $S$-\Term{morphism}{态射}。设 $x\in X$，且 $X$ 在 $x$ 处在 $S$ 上\Term{smooth}{光滑}。$g$ 在 $x$ 处\Term{etale}{平展}的充要条件是，各 $dg_i$（$1\leq i\leq n$）在 $x$ 处构成 $\Omega^1_{X/S}$ 的一组基；等价地，它们在
\[
\Omega^1_{X/S}(x)=(\Omega^1_{X/S})_x\otimes_{\mathcal O_x}k(x)
\]
中的像构成这个 $k(x)$-向量空间的一组基。
\end{SGAresult}

设 $X$ 是\Term{prescheme}{预概形}，$Y$ 是 $X$ 的闭子\Term{prescheme}{预概形}，由一个\Term{coherent}{凝聚}理想\Term{sheaf}{层} $\mathcal J$ 定义。于是 $\mathcal J/\mathcal J^2$ 可视为 $Y$ 上的\Term{coherent}{凝聚}\Term{sheaf}{层}，即 $Y$ 在 $X$ 中的\Term{conormal-sheaf}{余法层}。若 $X$ 还是 $S$-\Term{prescheme}{预概形}，则 $Y$ 上有如下典范的\Term{quasi-coherent}{拟凝聚}\Term{sheaf}{层}\Term{exact-sequence}{正合列}：
\[
\mathcal J/\mathcal J^2\xrightarrow{\ d\ }
\Omega^1_{X/S}\otimes_{\mathcal O_X}\mathcal O_Y
\longrightarrow\Omega^1_{Y/S}\longrightarrow0.\tag{4.3}
\]
右边这一段正是 (4.2)：交换 $X$、$Y$ 的角色，并注意 $\Omega^1_{Y/X}=0$。而\Term{homomorphism}{同态} $\mathcal J/\mathcal J^2\to\Omega^1_{X/S}\otimes_{\mathcal O_X}\mathcal O_Y$，则由通常并不线性的\Term{homomorphism}{同态} $g\mapsto dg$ 取商得到。(4.3) 的\Term{exactness}{正合性}是经典而直接的；在\Term{affine}{仿射}情形，它表述为对应于核为 $J$ 的 $A$-代数满\Term{homomorphism}{同态} $B\to C$ 的\Term{exact-sequence}{正合列}
\[
J/J^2\longrightarrow\Omega^1_{B/A}\otimes_B C
\longrightarrow\Omega^1_{C/A}\longrightarrow0,
\qquad C=B/J.\tag{4.3bis}
\]
这一\Term{exact-sequence}{正合列}其实已在 I 7.2 的证明中隐含地使用过！

\begin{SGAresult}{命题}{4.9}
设 $X$ 是 $S$-\Term{prescheme}{预概形}，$Y$ 是由 $X$ 上的\Term{coherent}{凝聚}理想\Term{sheaf}{层} $\mathcal J$ 定义的闭子\Term{prescheme}{预概形}；$x\in X$，$g_i$（$1\leq i\leq n$）是 $\mathcal O_X$ 的\Term{section}{截面}，定义 $S$-\Term{morphism}{态射}
\[
g:X\longrightarrow S[t_1,\ldots,t_n]=X';
\]
再设 $p$ 为整数，$0\leq p\leq n$。假设 $X$ 在 $x$ 处在 $S$ 上\Term{smooth}{光滑}。下列条件等价：
\begin{enumerate}
\item[(i)] $x$ 在 $X$ 中有一个\Term{open-neighborhood}{开邻域} $X_1$，使 $g|_{X_1}$ \Term{etale}{平展}，并且 $Y_1=Y\cap X_1$（$Y$ 在 $X_1$ 上的迹）为闭子\Term{prescheme}{预概形} $Y'=S[t_{p+1},\ldots,t_n]$ 在 $X'=S[t_1,\ldots,t_n]$ 中的\Term{inverse-image}{逆像}，即 $g_i$（$1\leq i\leq p$）生成 $\mathcal J|_{X_1}$：
\SourcePage{50}
\[
\begin{tikzcd}
Y_1 \arrow[r]\arrow[d]&X_1 \arrow[d,"\text{平展}"]\\
Y'=S[t_{p+1},\ldots,t_n]\arrow[r]&X'=S[t_1,\ldots,t_n].
\end{tikzcd}
\]
\item[(ii)] $Y$ 在 $x$ 处在 $S$ 上\Term{smooth}{光滑}；$g_i$（$1\leq i\leq p$）在 $\mathcal J_x$ 中定义元素；$dg_i(x)$（$1\leq i\leq n$）构成 $\Omega^1_{X/S}(x)$ 的一组 $k(x)$-基；$dg'_i(x)$（$p+1\leq i\leq n$）构成 $\Omega^1_{Y/S}(x)$ 的一组 $k(x)$-基。这里 $g'_i$ 表示 $g_i$ 在 $Y$ 上的限制，所有微分均相对于 $S$。
\item[(iii)] $g_i$（$1\leq i\leq p$）构成 $\mathcal J_x$ 的一组生成元，且 $dg_i(x)$（$1\leq i\leq n$）构成 $\Omega^1_{X/S}(x)$ 的一组 $k(x)$-基。
\item[(iv)] $Y$ 在 $x$ 处在 $S$ 上\Term{smooth}{光滑}；$g_i$（$1\leq i\leq p$）构成 $\mathcal J_x$ 的一组\Term{minimal-generating-system}{极小生成元系}；$dg'_i(x)$（$p+1\leq i\leq n$）构成 $\Omega^1_{Y/S}(x)$ 的一组 $k(x)$-基。
\end{enumerate}
此外，在这些条件下，$\mathcal J/\mathcal J^2$ 作为 $Y$ 上的模在 $x$ 处自由，$g_i$（$1\leq i\leq p$）的类构成它在 $x$ 处的一组基；典范\Term{homomorphism}{同态} $\mathcal J/\mathcal J^2\to\Omega^1_{X/S}\otimes\mathcal O_Y$ 在 $x$ 处\Term{universally-injective}{普遍单射}。
\end{SGAresult}

\textbf{注记。}这表明，$p$ 由其余条件唯一确定。它既是 $\mathcal J/\mathcal J^2$ 作为 $Y$ 上\Term{free-module}{自由模}在 $x$ 处的\Term{rank}{秩}，也是 $\mathcal J_x$ 在 $X$ 上所需生成元的最少数目；也可以用 $Y$ 相对于 $S$ 在 $x$ 处的\Term{relative-dimension}{相对维数}为 $n-p$ 来确定。

\begin{proof}
先设 (i) 成立。由 I 4.6 (iii)，$Y_1$ 在 $Y'$ 上\Term{etale}{平展}，故按定义，它在 $x$ 处在 $S$ 上\Term{smooth}{光滑}，\Term{relative-dimension}{相对维数}为 $n-p$；$Y$ 也如此。由 4.8，$dg_i$（$1\leq i\leq n$）在 $x$ 处构成 $\Omega^1_{X/S}$ 的一组基，$dg'_i$（$p+1\leq i\leq n$）在 $x$ 处构成 $\Omega^1_{Y/S}$ 的一组基。于是由\Term{exact-sequence}{正合列} (4.3)，$g_i$（$1\leq i\leq p$）在视为 $Y$-模的 $\mathcal J/\mathcal J^2$ 中在 $x$ 处线性无关。又因它们生成 $\mathcal J_x$，可知各 $g_i\bmod\mathcal J_x^2$ 构成 $\mathcal J/\mathcal J^2$ 在 $x$ 处的一组基。一方面，这说明 $g_i$（$1\leq i\leq p$）构成 $\mathcal J_x$ 的\Term{minimal-generating-system}{极小生成元系}；另一方面，(4.3) 中的\Term{homomorphism}{同态} $\mathcal J/\mathcal J^2\to\Omega^1_{X/S}\otimes\mathcal O_Y$ 在 $x$ 处\Term{universally-injective}{普遍单射}，因为它把一个在 $x$ 处自由的模的一组基映到另一个在 $x$ 处自由的模的一组基的一部分。注意，这里说的都是 $Y$-模。于是 (i) 蕴含 (ii)、(iii)、(iv)，并且蕴含命题 4.9 最后的各项断言。

由推论 4.8，(iii) 蕴含 (i)。

(ii) 蕴含 (i)：(ii) 中的第一个假设意味着，必要时以 $x$ 在 $X$ 中的某个\Term{open-neighborhood}{开邻域}替换 $X$ 后，$g$ 诱导\Term{morphism}{态射} $h:Y\to Y'$。由 4.8，(ii) 中其余两个假设意味着 $g$ 和 $h$ 都在 $x$ 处\Term{etale}{平展}。令 $Y''$ 为 $Y'$ 在 $g$ 下的\Term{inverse-image}{逆像}。则 $Y$ 是 $Y''$ 的闭子\Term{prescheme}{预概形}；因为 $g$ 在 $x$ 处\Term{etale}{平展}，由 I 4.6 (iii)，$Y''$ 在 $x$ 处在 $Y'$ 上\Term{etale}{平展}。因此\Term{immersion}{浸入}\Term{morphism}{态射} $Y\to Y''$ 本身\Term{etale}{平展}（I 4.8），从而是
\SourcePage{51}
\Term{open-immersion}{开浸入}（I 5.8 或 I 5.2）。再以 $x$ 的某个合适\Term{open-neighborhood}{开邻域} $X_1$ 替换 $X$，便得到 (i)。

以上已证明 (i)、(ii)、(iii) 等价，并且都蕴含 (iv)。只须再证 (iv)$\Rightarrow$(ii)。考虑到 $\Omega^1_{X/S}$ 作为 $X$ 上的模在 $x$ 处自由，只要知道“$Y$ 在 $x$ 处在 $S$ 上\Term{smooth}{光滑}”蕴含“$\mathcal J/\mathcal J^2$ 作为 $Y$ 上的模在 $x$ 处自由，且 $\mathcal J/\mathcal J^2\to\Omega^1_{X/S}\otimes\mathcal O_Y$ 在 $x$ 处\Term{universally-injective}{普遍单射}”，结论立即成立。这最后一点包含在下面的定理中。
\end{proof}

\begin{SGAresult}{定理}{4.10}
设 $X$ 是\Term{smooth}{光滑} $S$-\Term{prescheme}{预概形}，$Y$ 是由 $X$ 上的\Term{coherent}{凝聚}理想\Term{sheaf}{层} $\mathcal J$ 定义的闭子\Term{prescheme}{预概形}，$x\in X$。下列条件等价：
\begin{enumerate}
\item[(i)] $Y$ 在 $x$ 处在 $S$ 上\Term{smooth}{光滑}。
\item[(ii)] $x$ 在 $X$ 中有一个\Term{open-neighborhood}{开邻域} $X_1$，并有\Term{etale}{平展} $S$-\Term{morphism}{态射}
\[
g:X_1\longrightarrow X'=S[t_1,\ldots,t_n],
\]
使 $Y_1=Y\cap X_1$（$Y$ 在 $X_1$ 上的迹）为 $X_1$ 的子\Term{prescheme}{预概形}，它等于某个闭子\Term{prescheme}{预概形} $Y'=S[t_{p+1},\ldots,t_n]$ 在 $X'=S[t_1,\ldots,t_n]$ 中经 $g$ 的\Term{inverse-image}{逆像}，其中 $p$ 取适当值。
\item[(iii)] $\mathcal J_x$ 有生成元 $g_i$（$1\leq i\leq p$），使各 $dg_i$ 在 $x$ 处构成 $\Omega^1_{X/S}$ 的一组基的一部分；等价地，各 $dg_i(x)$ 在 $\Omega^1_{X/S}$ 中在 $k(x)$ 上线性无关。
\item[(iv)] $\mathcal J/\mathcal J^2$ 作为 $Y$ 上的\Term{sheaf}{层}在 $x$ 处自由，且典范\Term{homomorphism}{同态}
\[
d:\mathcal J/\mathcal J^2\longrightarrow\Omega^1_{X/S}\otimes\mathcal O_Y
\]
在 $x$ 处\Term{universally-injective}{普遍单射}。等价地，典范\Term{homomorphism}{同态}列
\[
0\longrightarrow\mathcal J/\mathcal J^2
\longrightarrow\Omega^1_{X/S}\otimes\mathcal O_Y
\longrightarrow\Omega^1_{Y/S}\longrightarrow0
\]
在 $x$ 处\Term{exact}{正合}，且 $\Omega^1_{Y/S}$ 在 $x$ 处\Term{locally-free}{局部自由}。
\end{enumerate}
\end{SGAresult}
\begin{proof}
由上面的论证，已知 (ii) 蕴含 (i)、(iii)、(iv)。现证 (i)$\Rightarrow$(ii)，这也同时完成 4.9 的证明。由定理 4.3 (ii)，\Term{exact-sequence}{正合列} (4.3) 中最后两项是 $Y$ 上的\Term{free-module}{自由模}。各 $dg$（$g\in\mathcal O_X$）在 $\Omega^1_{X/S}\otimes_{\mathcal O_X}\mathcal O_Y$ 中的像在 $x$ 处生成这个模，故它们在 $\Omega^1_{Y/S}$ 中的像也在 $x$ 处生成后者。因此可取 $\mathcal O_X$ 中的 $g_i$（$p+1\leq i\leq n$），使 $dg'_i$ 构成 $\Omega^1_{Y/S}$ 的一组基，再由 (4.3) 的\Term{exactness}{正合性}，用形如 $dg_i$（$1\leq i\leq n$）、且 $g_i\in\mathcal J_x$（$1\leq i\leq p$）的元素，把 $dg_i$（$p+1\leq i\leq n$）补成中间一项的一组基。各 $g_i$ 来自 $\mathcal O_X$ 在 $x$ 的某个邻域上的\Term{section}{截面}；可设该邻域就是 $X$。于是满足 4.8 (ii) 的条件，而我们已证明这蕴含 4.8 (i)，从而得到 4.10 (ii)。
\SourcePage{52}
4.10 中的 (iii)$\Rightarrow$(ii) 立即来自 4.8 中的 (iii)$\Rightarrow$(i)。因此 (i)、(ii)、(iii) 等价，并且蕴含 (iv)。最后，(iv) 显然蕴含 (iii)，只须注意：$g_i\in\mathcal J_x$ 若模 $\mathcal J_x^2$ 后构成一组基，就生成 $\mathcal J_x$（Nakayama 引理）。
\end{proof}

此外，上面的证明还表明：
\begin{SGAresult}{推论}{4.11}
设 $X$ 是 $S$-\Term{prescheme}{预概形}，$Y$ 是由 $X$ 上的\Term{coherent}{凝聚}理想\Term{sheaf}{层} $\mathcal J$ 定义的闭子\Term{prescheme}{预概形}，$x\in Y$。设 $X$、$Y$ 都在 $x$ 处在 $S$ 上\Term{smooth}{光滑}。令 $g_i$（$1\leq i\leq p$）为 $\mathcal J$ 的\Term{section}{截面}。下列条件等价：
\begin{enumerate}
\item[(i)] 各 $g_i$ 生成 $\mathcal J_x$，且各 $dg_i(x)$ 在 $\Omega^1_{X/S}(x)$ 中在 $k(x)$ 上线性无关。
\item[(ii)] 各 $g_i\bmod\mathcal J^2$ 在 $x$ 处构成 $\mathcal J/\mathcal J^2$ 的一组基。
\item[(iii)] 各 $g_i$ 构成 $\mathcal J_x$ 的一组\Term{minimal-generating-system}{极小生成元系}。
\item[(iv)] 可以在 $X$ 的某个邻域 $X_1$ 上另取 $\mathcal O_X$ 的\Term{section}{截面} $g_i$（$p+1\leq i\leq n$），使其与已有的\Term{section}{截面}一起定义一个\Term{etale}{平展}\Term{morphism}{态射} $X_1\to X'=S[t_1,\ldots,t_n]$，且 $Y_1=Y\cap X_1$ 是闭子\Term{prescheme}{预概形} $Y'=S[t_{p+1},\ldots,t_n]$ 在 $X'=S[t_1,\ldots,t_n]$ 中经 $g$ 的\Term{inverse-image}{逆像}。
\end{enumerate}
\end{SGAresult}

特别地，有：
\begin{SGAresult}{推论}{4.12}
设 $X$ 是 $S$-\Term{prescheme}{预概形}，$F$ 是 $\mathcal O_X$ 的\Term{section}{截面}，$Y$ 是 $F$ 的零点子\Term{prescheme}{预概形}，即由\Term{coherent}{凝聚}理想 $F\mathcal O_X$ 定义的子\Term{prescheme}{预概形}，$x\in Y$。设 $X$ 在 $x$ 处在 $S$ 上简单。则 $Y$ 在 $x$ 处在 $S$ 上\Term{smooth}{光滑}的充要条件是：$F$ 在 $x$ 的某个邻域上为零，或 $dF(x)\ne0$；这里 $dF(x)$ 指 $dF$ 在 $k(x)$-向量空间 $\Omega^1_{X/S}(x)$ 中的像。
\end{SGAresult}
充分性来自 4.10 的准则 (iii)。证明必要性。因为 $\mathcal J$ 由一个元素生成，$\mathcal J/\mathcal J^2$ 首先必须在 $x$ 处自由，且\Term{rank}{秩}不超过 $1$。若\Term{rank}{秩}为 $0$，即 $\mathcal J/\mathcal J^2$ 在 $x$ 处为零，则由 Nakayama 引理，$\mathcal J$ 在 $x$ 处为零，即 $F$ 在 $x$ 的某个邻域上为零。若\Term{rank}{秩}为 $1$，则 $F$ 构成 $\mathcal J$ 在 $x$ 处的\Term{minimal-generating-system}{极小生成元系}，再由 4.11 中 (i)、(iii) 的等价性得出结论。

\begin{SGAresult}{推论}{4.13}
设 $Y$ 是\Term{locally-finite-type}{局部有限型} $S$-\Term{prescheme}{预概形}，$S'$ 是\Term{flat}{平坦} $S$-\Term{prescheme}{预概形}，$Y'=Y\times_S S'$；$x'\in Y'$，$x$ 是它在 $Y$ 中的典范像。$Y$ 在 $x$ 处在 $S$ 上\Term{smooth}{光滑}的充要条件是 $Y'$ 在 $x'$ 处在 $S'$ 上\Term{smooth}{光滑}。特别地，若 $S'\to S$ \Term{flat}{平坦}且\Term{surjectivity}{满射}，则 $Y$ 在 $S$ 上\Term{smooth}{光滑}当且仅当 $Y'$ 在 $S'$ 上\Term{smooth}{光滑}。
\end{SGAresult}
只须证明充分性；必要性已见于 1.3 (ii)。以 $x$ 的某个合适邻域替换 $Y$，并以其\Term{inverse-image}{逆像}替换 $Y'$，可设 $S$ \Term{affine}{仿射}，$Y$ 在 $S$ 上\Term{affine}{仿射}且\Term{finite-type}{有限型}。于是 $Y$ \Term{isomorphism}{同构}于某个形如 $S[t_1,\ldots,t_n]$ 的\Term{scheme}{概形}的闭子\Term{prescheme}{预概形}。
\SourcePage{53}
因而 $Y'$ 等同于 $X'=X\times_S S'$ 的一个闭子\Term{prescheme}{预概形}。$X$ 在 $S$ 上\Term{smooth}{光滑}，故 $X'$ 在 $S'$ 上\Term{smooth}{光滑}，可以应用\Term{smoothness}{光滑性}准则 4.10；这里准则 (iv) 立即给出结果。

\textbf{注记 4.14。}定理 4.10 的准则 (iii) 值得称为\Term{jacobian-criterion}{光滑性的雅可比准则}。从理论上说，它可以判定给定 $S$-\Term{prescheme}{预概形} $Y$ 是否在点 $x\in Y$ 处在 $S$ 上\Term{smooth}{光滑}，因为 $Y$ 总有一个邻域\Term{isomorphism}{同构}于某个\Term{smooth}{光滑} $S$-\Term{prescheme}{预概形} $X$ 的子\Term{prescheme}{预概形}，例如 $X=S[t_1,\ldots,t_n]$。人们通常也正是在 $X=S[t_1,\ldots,t_n]$、$S=\operatorname{Spec}(A)$ 时叙述\Term{jacobian-criterion}{雅可比准则}；当然，Zariski 所考虑的经典情形中，$A$ 是域。给定 $A[t_1,\ldots,t_n]$ 的一个理想 $J$ 及包含 $J$ 的一个\Term{prime-ideal}{素理想}后，相应的表述留给读者写出。还应指出，就目前看来，特别是在 Nagata 已用非微分方法推广了 Zariski 关于代数\Term{scheme}{概形}的\Term{regular-point}{正则点}集为开集的定理之后，\Term{jacobian-criterion}{雅可比准则}几乎只有在这里给出的形式下才有意义：仅使用\Term{relative-differential}{相对微分}，而不使用\Term{absolute-differential}{绝对微分}，即相对于绝对常数环 $\mathbb Z$ 的微分。和许多情形一样，这里考虑微分比考虑\Term{derivation}{导子}更方便。最后注意，若 $Y$ 在 $x$ 处在 $S$ 上\Term{smooth}{光滑}，\Term{relative-dimension}{相对维数}为 $n$，则 $x$ 在 $Y$ 中有\Term{open-neighborhood}{开邻域}\Term{isomorphism}{同构}于 $X=S[t_1,\ldots,t_n]$ 的一个子\Term{prescheme}{预概形}，其中 $n=m+1$；这由定义和 I 7.6 得出。

设 $A$ 是\Term{noetherian-ring}{诺特环}，$x_i$（$1\leq i\leq n$）是 $A$ 中的元素，$J$ 是它们生成的理想。若由各 $x_i$ 定义的典范满\Term{homomorphism}{同态}
\[
(A/J)[t_1,\ldots,t_n]\longrightarrow\operatorname{gr}_J(A)
\]
是\Term{isomorphism}{同构}，则称各 $x_i$ 构成 $J$ 的一组\Term{regular-generating-system}{正则生成元系}。右端是 $A$ 关于 $J$ 的幂所给滤过的\Term{associated-graded-ring}{伴随分次环}。这一条件也就是说：
\begin{enumerate}
\item[(i)] 典范满\Term{homomorphism}{同态}
\[
\operatorname{S}_{A/J}(J/J^2)\longrightarrow\operatorname{gr}_J(A)
\]
是\Term{isomorphism}{同构}；左端表示 $A/J$-模 $J/J^2$ 的\Term{symmetric-algebra}{对称代数}；
\item[(ii)] $J/J^2$ 自由，且各 $x_i\bmod J^2$ 为其一组基。
\end{enumerate}
由这种表述可见，若 $J\ne A$，则各 $x_i$ 构成 $J$ 的一组\Term{minimal-generating-system}{极小生成元系}，而 $J$ 的任意另一组\Term{minimal-generating-system}{极小生成元系}也都是\Term{regular-generating-system}{正则生成元系}。注意，这里的“极小”取严格意义，即元素数目最少；只有当 $A$ 是\Term{local-ring}{局部环}时，它才等价于按包含关系的极小性。另一方面，若 $J=A$，则 $J$ 的任意一组生成元都是\Term{regular}{正则}的。

\SourcePage{54}
理想的一组生成元为\Term{regular-generating-system}{正则生成元系}这一条件，在对任意乘法封闭集合作\Term{localization}{局部化}后保持。另一方面，立即可见，要使 $(x_i)$ 是 $J$ 的一组\Term{minimal-generating-system}{极小生成元系}，只须对每个包含 $J$ 的\Term{maximal-ideal}{极大理想} $\mathfrak m$，各 $x_i$ 在 $A_{\mathfrak m}$ 中构成 $JA_{\mathfrak m}$ 的一组\Term{regular-generating-system}{正则生成元系}。因此可归结为 $A$ 是\Term{maximal-ideal}{极大理想}为 $\mathfrak m$ 的\Term{local-ring}{局部环}，且各 $x_i$ 都属于 $\mathfrak m$ 的情形。此时，各 $x_i$ 构成 $J$ 的一组\Term{regular-generating-system}{正则生成元系}，当且仅当它们构成 Serre 意义下的 $A$-序列\footnote{现在更常说“$A$-\Term{regular-sequence}{正则序列}”，见 EGA $0_{\mathrm{IV}}$ 15.1.7 与 15.1.11。}，即对每个 $1\leq i\leq n$，$x_i$ 在 $A/(x_1,\ldots,x_{i-1})A$ 中不是\Term{zero-divisor}{零因子}。

最后，若 $A$ 是环 $B$ 上的代数，且 $A/J$ 作为 $B$-代数\Term{isomorphism}{同构}于 $B$，从而 $J$ 是某个 $B$-代数\Term{homomorphism}{同态} $A\to B$ 的核，则各 $x_i$ 构成 $J$ 的一组\Term{regular-generating-system}{正则生成元系}，当且仅当由它们定义的典范\Term{homomorphism}{同态}
\[
B[[t_1,\ldots,t_n]]\longrightarrow\widehat A
\]
是\Term{isomorphism}{同构}；这个\Term{homomorphism}{同态}无论如何总是\Term{surjectivity}{满射}。右端是 $A$ 对 $J$ 的幂所定义的拓扑的\Term{separated}{分离}\Term{completion}{完备化} $\varprojlim A/J^{n+1}$。

这些事实都很熟悉，在 Gabriel 整理的 Serre 交换代数讲义中应该大体都有；在 $A$ 为\Term{local-ring}{局部环}时，那份讲义还给出了 $A$-序列的另外许多刻画。

设 $J$ 是\Term{noetherian-ring}{诺特环} $A$ 的理想。若对 $A$ 的每个\Term{prime-ideal}{素理想} $\mathfrak p$，$JA_{\mathfrak p}$ 都有\Term{regular-generating-system}{正则生成元系}，则称 $J$ 是\Term{regular-ideal}{正则理想}。显然，只须对 $\mathfrak p\supset J$ 核对，而且还可限于\Term{maximal-ideal}{极大理想}。更一般地，若 $\mathcal J$ 是\Term{locally-noetherian}{局部诺特}\Term{prescheme}{预概形} $X$ 上的理想\Term{sheaf}{层}，则称 $\mathcal J$ 为\Term{regular-ideal}{正则理想}，是指对每个 $x\in X$，$\mathcal J_x$ 都是 $\mathcal O_x$ 的一个具有\Term{regular-generating-system}{正则生成元系}的理想。这等价于以下两个条件同时成立：
\begin{enumerate}
\item[(a)] 典范满\Term{homomorphism}{同态}
\[
\operatorname{S}_{\mathcal O_X/\mathcal J}(\mathcal J/\mathcal J^2)
\longrightarrow\operatorname{gr}_{\mathcal J}(\mathcal O_X)
\]
是\Term{isomorphism}{同构}；
\item[(b)] $\mathcal O_X/\mathcal J$-模\Term{sheaf}{层} $\mathcal J/\mathcal J^2$ \Term{locally-free}{局部自由}。
\end{enumerate}
此时也称由 $\mathcal J$ 定义的 $X$ 的子\Term{prescheme}{预概形} $Y$ \Term{regularly-immersed}{正则浸入}于 $X$；这里 $\mathcal O_Y$ 以零延拓后\Term{isomorphism}{同构}于 $\mathcal O_X/\mathcal J$。同样可显然地定义\Term{regular-immersion}{正则浸入}\Term{morphism}{态射}，以及“在点 $x$ 处\Term{regular}{正则}”：\Term{immersion}{浸入} $Y\to X$ 把 $Y$（分别地，$x$ 的某个合适邻域）等同于 $X$ 的某个开集中的一个\Term{regularly-immersed}{正则浸入}的闭子\Term{prescheme}{预概形}。这里不能说“\Term{regular}{正则}子\Term{prescheme}{预概形}”，因为那是指 $Y$ 的\Term{local-ring}{局部环}\Term{regular}{正则}。最后，若对每个 $x\in X$，$\mathcal J$ 的\Term{section}{截面} $x_i$ 在 $\mathcal O_x$ 中的相应元素构成 $\mathcal J_x$ 的一组\Term{regular-generating-system}{正则生成元系}，则称这些\Term{section}{截面}构成\Term{regular-generating-system}{正则生成元系}。

\SourcePage{55}
这与对环的理想生成元所引入的术语一致。它也就是说，由各 $x_i$ 定义的典范满\Term{homomorphism}{同态}
\[
\mathcal O_Y[t_1,\ldots,t_n]\longrightarrow\operatorname{gr}_{\mathcal J}(\mathcal O_X)
\]
是\Term{isomorphism}{同构}。若预先知道理想\Term{sheaf}{层} $\mathcal J$ \Term{regular}{正则}，那么这只意味着：在 $Y$ 的每一点 $x$，各 $x_i$ 给出 $\mathcal J/\mathcal J^2$ 在 $\mathcal O_{Y,x}$ 上的一组基。注意，若 $Y$ 为空，这个条件是空的。因此，$\mathcal J$ 有\Term{regular-generating-system}{正则生成元系}的充要条件是：$\mathcal J$ \Term{regular}{正则}，且 $\mathcal O_Y$-模 $\mathcal J/\mathcal J^2$ 整体自由，而非仅\Term{locally-free}{局部自由}；也就是说，典范\Term{homomorphism}{同态} $\operatorname{S}_{\mathcal O_Y}(\mathcal J/\mathcal J^2)\to\operatorname{gr}_{\mathcal J}(\mathcal O_X)$ 为\Term{surjectivity}{满射}，且 $\mathcal O_Y$-模 $\mathcal J/\mathcal J^2$ 整体自由。

一个\Term{augmented-ring}{增广环}称为\Term{regular}{正则}的，是指它的\Term{augmentation-ideal}{增广理想}\Term{regular}{正则}。例如，\Term{local-ring}{局部环} $A$ 配以到其\Term{residue-field}{剩余域} $k$ 的增广后，$A$ 是\Term{regular-local-ring}{正则局部环}当且仅当它是\Term{regular}{正则}\Term{augmented-ring}{增广环}。

其实，似乎没有必要先对环罗列这一整套预备性质；直接从\Term{sheaf}{层}开始更为合适。在\Term{noetherian}{诺特}情形，如果要有一个适于微分计算的定义，这里采用的定义——先验上弱于用 Serre 的 $A$-序列给出的定义——似乎更好。当然，若想把事情做好，还应至少把\Term{smooth}{光滑}\Term{morphism}{态射}理论的一部分推广到非\Term{noetherian}{诺特}情形\footnote{正如前言所说，这项工作现在已完成，见 EGA IV 17、18。}；很可能应从\Term{jacobian-criterion}{雅可比准则}出发，争取获得\Term{smooth}{光滑}\Term{morphism}{态射}和\Term{etale}{平展}\Term{morphism}{态射}，即\Term{smooth}{光滑}且\Term{quasi-finite}{拟有限}的\Term{morphism}{态射}，所有主要的形式性质，让\Term{noetherian}{诺特}假设只在逆命题中出现。

经过这番冗长的术语准备，下面给出一个小定理。
\begin{SGAresult}{定理}{4.15}
设 $X$ 是\Term{locally-finite-type}{局部有限型} $S$-\Term{prescheme}{预概形}，$Y$ 是由 $X$ 上的\Term{coherent}{凝聚}理想\Term{sheaf}{层} $\mathcal J$ 定义的闭子\Term{prescheme}{预概形}，$x\in X$。现在假设 $Y$ 在 $x$ 处在 $S$ 上\Term{smooth}{光滑}，而不对 $X$ 作任何假设。则下列条件等价：
\begin{enumerate}
\item[(i)] $X$ 在 $x$ 处在 $S$ 上\Term{smooth}{光滑}；
\item[(ii)] \Term{immersion}{浸入} $i:Y\to X$ 在 $x$ 处\Term{regular}{正则}，即 $\mathcal J_x$ 是 $\mathcal O_x$ 的\Term{regular-ideal}{正则理想}。
\end{enumerate}
\end{SGAresult}

\begin{SGAresult}{推论}{4.16}
设 $Y$ 在 $S$ 上\Term{smooth}{光滑}。$X$ 在 $Y$ 的某个邻域中在 $S$ 上\Term{smooth}{光滑}，即在 $Y$ 的各点都\Term{smooth}{光滑}，当且仅当 $Y$ \Term{regularly-immersed}{正则浸入}于 $X$，即\Term{immersion}{浸入} $i:Y\to X$ \Term{regular}{正则}。
\end{SGAresult}

\SourcePage{56}
\begin{proof}
(i) 蕴含 (ii)。应用 4.10 的准则 (ii)。因为 $g:X_1\to X$ \Term{flat}{平坦}，要证明 $X'$ 的子\Term{prescheme}{预概形} $Y'$ 经 $g$ 的\Term{inverse-image}{逆像}\Term{regularly-immersed}{正则浸入}，可归结为证明 $Y'=S[t_{p+1},\ldots,t_n]$ \Term{regularly-immersed}{正则浸入}于 $S[t_1,\ldots,t_n]$。这显然成立：各 $t_i$（$1\leq i\leq p$）构成在 $X'$ 中定义 $Y'$ 的理想\Term{sheaf}{层}的一组\Term{regular-generating-system}{正则生成元系}。

(ii) 蕴含 (i)。取 $\mathcal J_x$ 的一组\Term{regular-generating-system}{正则生成元系} $g_i$（$1\leq i\leq p$），再取 $\mathcal O_{X,x}$ 中的元素 $g_i$（$p+1\leq i\leq n$），使它们在 $\mathcal O_{Y,x}$ 中的像 $g'_i$ 定义一个\Term{etale}{平展}\Term{morphism}{态射}
\[
Y_1\longrightarrow Y'=S[t_{p+1},\ldots,t_n],
\]
从 $Y$ 的一个邻域 $Y_1$ 映到 $Y'$。各 $g_i$（$1\leq i\leq n$）来自 $\mathcal O_X$ 在 $x$ 的某个邻域 $X_1$ 上的同名\Term{section}{截面}，可设 $X_1=X$、$Y_1=Y$。于是得到\Term{morphism}{态射}
\[
g:X\longrightarrow X'=S[t_1,\ldots,t_n],
\]
只须证明它在 $x$ 处\Term{etale}{平展}。把 $X_1$ 取得足够小，可设各 $g_i$（$1\leq i\leq p$）在整个 $X$ 上构成 $\mathcal J$ 的一组\Term{regular-generating-system}{正则生成元系}。特别地，它们生成 $\mathcal J$，因此 $X$ 的子\Term{prescheme}{预概形} $Y$ 等同于 $X'$ 的子\Term{prescheme}{预概形} $Y'$ 经 $g$ 的\Term{inverse-image}{逆像}。令 $x'=g(x)$，则 $X'\to X$ 在 $x'$ 处的\Term{fiber}{纤维}与 $Y\to Y'$ 在 $x$ 处的\Term{fiber}{纤维}相同，因而在 $k(x')$ 上\Term{etale}{平展}。所以 $g$ 在 $x$ 处\Term{unramified}{非分歧}，只须再证它在 $x$ 处\Term{flat}{平坦}。

$\mathcal O_{X',x'}$ 关于 $\mathcal J'_{x'}$ 的幂的\Term{associated-graded-ring}{伴随分次环}，在每一度都是自由 $\mathcal O_{Y',x'}$-模。另一方面，$\mathcal O_{X,x}$ 关于 $\mathcal J_x=\mathcal J'_{x'}\mathcal O_{X,x}$ 的幂的\Term{associated-graded-ring}{伴随分次环}，在典范\Term{homomorphism}{同态}下\Term{isomorphism}{同构}于前一个\Term{associated-graded-ring}{伴随分次环}与 $\mathcal O_{Y,x}$ 的\Term{tensor-product}{张量积}：这两个环都是有 $n-p$ 个不定元的多项式环，常数环分别为 $\mathcal O_{Y',x'}$ 和 $\mathcal O_{Y,x}$。最后，$\mathcal O_{X,x}/\mathcal J_x=\mathcal O_{Y,x}$ 在 $\mathcal O_{X',x'}/\mathcal J'_{x'}=\mathcal O_{Y',x'}$ 上\Term{flat}{平坦}。

由一个一般的\Term{flatness}{平坦性}准则，可知 $X$ 在 $x$ 处在 $X'$ 上\Term{flat}{平坦}，证毕。这个准则适用于\Term{noetherian-local-ring}{诺特局部环}的\Term{local-homomorphism}{局部同态} $A'\to A$，其中 $A'$ 带有理想 $J'\ne A'$，且相应\Term{associated-graded-ring}{伴随分次环}在每一度都是自由 $A'/J'$-模。
\end{proof}

\begin{SGAresult}{推论}{4.17}
设 $X$ 是 $Y$ 上\Term{locally-finite-type}{局部有限型}的\Term{prescheme}{预概形}，$i$ 是 $X$ 在 $Y$ 上的一个\Term{section}{截面}，$y\in Y$，$x=i(y)$，$\mathcal J$ 是由子\Term{prescheme}{预概形} $i(Y)$ 定义的 $X$ 上的理想\Term{sheaf}{层}。为简化表述，假设 $i(Y)$ 闭；当 $X$ 是\Term{scheme}{概形}时，该条件成立。下列条件等价：
\begin{enumerate}
\item[(i)] $X$ 在 $x$ 处在 $Y$ 上\Term{smooth}{光滑}。
\item[(ii)] $i$ 在 $y$ 处是\Term{regular-immersion}{正则浸入}。
\item[(iii)] $\mathcal O_x$ 关于 $\mathcal J_x$ 的幂所定义拓扑的\Term{completion}{完备化}，作为 $\mathcal O_y$-代数\Term{isomorphism}{同构}于一个\Term{formal-power-series-algebra}{形式幂级数代数} $\mathcal O_y[[t_1,\ldots,t_n]]$。
\SourcePage{57}
\item[(iii bis)] $y$ 有一个\Term{open-neighborhood}{开邻域} $U$，使 $\mathcal O_Y$ 上的代数\Term{sheaf}{层} $\varprojlim i^*(\mathcal O_X/\mathcal J^{n+1})$ 在 $U$ 上\Term{isomorphism}{同构}于形如 $\mathcal O_Y[[t_1,\ldots,t_n]]$ 的\Term{sheaf}{层}。
\item[(iv)] 存在 $y$ 的\Term{open-neighborhood}{开邻域} $U$、$x$ 的\Term{open-neighborhood}{开邻域} $V$，以及 $Y$-\Term{morphism}{态射} $g:V\to U[t_1,\ldots,t_n]$，使 $g$ \Term{etale}{平展}，$i$ 诱导 $V$ 在 $U$ 上的一个\Term{section}{截面}，且该\Term{section}{截面}经 $g$ 映成 $U[t_1,\ldots,t_n]$ 在 $U$ 上的零\Term{section}{截面}。
\end{enumerate}
\end{SGAresult}
(i)、(ii) 的等价性是定理 4.15 在 $Y=S$ 时的特例；(ii)、(iii) 的等价性，以及实质上 (ii)、(iii bis) 的等价性，已在上面的“回顾”中指出。(i)、(iv) 的等价性，则由定理 4.10 中 (i)、(ii) 的等价性容易得到。

\begin{SGAresult}{推论}{4.18}
设 $X$ 是 $S$ 上的\Term{smooth}{光滑}\Term{prescheme}{预概形}。则\Term{diagonal-morphism}{对角态射}
\[
\Delta_{X/S}:X\longrightarrow X\times_S X
\]
是\Term{regular-immersion}{正则浸入}；换一种说法，$X$ 在 $S$ 上“\Term{differentially-smooth}{微分光滑}”。
\end{SGAresult}
这确是推论 4.16 的特例，因为 $X$ 与 $X\times_S X$ 都在 $S$ 上\Term{smooth}{光滑}。

\textbf{注记 4.18。}回顾第 I 讲第 1 节：若 $X$ 是 $S$ 上的\Term{prescheme}{预概形}，则引入 $X$ 上的\Term{quasi-coherent}{拟凝聚}代数\Term{sheaf}{层}
\[
\mathcal P^n_{X/S}=\mathcal O_{X\times_S X}/\mathcal I_X^{n+1},
\]
其中 $\mathcal I_X$ 是在 $X\times_S X$ 中定义对角线的理想\Term{sheaf}{层}，并借助第一投影 $\operatorname{pr}_1:X\times_S X\to X$ 将它视为 $\mathcal O_X$-代数\Term{sheaf}{层}。各 $\mathcal P^n_{X/S}$ 构成 $X$ 上的代数\Term{projective-system}{逆系}，其\Term{projective-limit}{逆极限}记为 $\mathcal P^\infty_{X/S}$。它正是 $X\times_S X$ 沿对角线的\Term{formal-completion}{形式完备化}的结构\Term{sheaf}{层}；这里现在假定 $X$ 在 $S$ 上\Term{locally-finite-type}{局部有限型}，所以各 $\mathcal P^n_{X/S}$ 都\Term{coherent}{凝聚}。说 $X$ 在 $S$ 上\Term{differentially-smooth}{微分光滑}，即\Term{diagonal-morphism}{对角态射} $\Delta_{X/S}$ 为\Term{regular-immersion}{正则浸入}，也就是说，$\mathcal P^\infty_{X/S}$ 作为向 $\mathcal O_X$ 增广的代数\Term{sheaf}{层}是\Term{regular}{正则}的；也即 $\Omega^1_{X/S}$ \Term{locally-free}{局部自由}，且典范满\Term{homomorphism}{同态}
\[
\operatorname{S}_{\mathcal O_X}(\Omega^1_{X/S})
\longrightarrow\operatorname{gr}^*(\mathcal P^\infty_{X/S})
\]
是\Term{isomorphism}{同构}。还等价于说，$X$ 的每一点都有一个\Term{open-neighborhood}{开邻域}，使\Term{augmented-algebra}{增广代数}\Term{sheaf}{层} $\mathcal P^\infty_{X/S}$ 在其上\Term{isomorphism}{同构}于某个 $\mathcal O_X[[t_1,\ldots,t_n]]$。

设 $s$ 是 $X$ 在 $S$ 上的\Term{section}{截面}，$\mathcal J$ 是它定义的 $X$ 上的理想\Term{sheaf}{层}；为简化起见，假定 $s(S)$ 闭。则有如下典范的\Term{augmented-algebra}{增广} $\mathcal O_X$-代数\Term{isomorphism}{同构}：
\[
s^*(\mathcal P^n_{X/S})=\mathcal O_X/\mathcal J^{n+1},\qquad
s^*(\mathcal P^\infty_{X/S})=\varprojlim_n\mathcal O_X/\mathcal J^{n+1}.\tag{4.4}
\]
这些\Term{isomorphism}{同构}在\Term{base-change}{基变换}下具有显然意义上的\Term{functorial}{函子性}；结合这个事实，它们再次给出了 $S$ 上的代数\Term{sheaf}{层} $\mathcal P^n_{X/S}$ 的一种刻画。
\SourcePage{58}
例如，若 $S=\operatorname{Spec}(k)$，$k$ 为域，则给出 $X$ 在 $S$ 上的\Term{section}{截面} $s$，等价于给出 $X$ 的一个 $k$-有理点 $x$；上述公式意味着有 $k$-代数\Term{isomorphism}{同构}
\[
\mathcal P^n_{X/S}(x)=\mathcal O_x/\mathfrak m_x^{n+1}.\tag{4.5}
\]
这说明，把 $\mathcal P^n_{X/S}$ 称为“$X$ 相对于 $S$ 的 $n$ 阶\Term{principal-parts-sheaf}{主部层}”是恰当的。由 (4.4) 还可见，若 $X$ 在 $s(S)$ 的每一点都在 $S$ 上\Term{differentially-smooth}{微分光滑}，则 $X$ 在 $s(S)$ 的每一点都在 $S$ 上\Term{smooth}{光滑}（推论 4.17）；逆命题也成立（推论 4.18）。结合 4.13，不难得出：对于\Term{locally-finite-type}{局部有限型} $S$-\Term{prescheme}{预概形} $X$，$X$ 在 $S$ 上\Term{smooth}{光滑}当且仅当它在 $S$ 上\Term{flat}{平坦}且\Term{differentially-smooth}{微分光滑}。注意，\Term{flatness}{平坦性}假设不可缺少；取 $X$ 为 $S$ 的一个闭子\Term{prescheme}{预概形}即可看出这一点。

还要回顾，借助投影 $\operatorname{pr}_2:X\times_S X\to X$，可在 $\mathcal P^n_{X/S}$ 上得到第二种代数结构。事实上，它可由第一种结构通过环\Term{sheaf}{层} $\mathcal P^n_{X/S}$ 的典范对合得到；这个对合由 $X\times_S X$ 的交换两因子的\Term{automorphism}{自同构}诱导。以 $d^n_{X/S}$ 或简记 $d^n$ 表示对应于第二种代数结构的环\Term{sheaf}{层}\Term{homomorphism}{同态}
\[
d^n_{X/S}:\mathcal O_X\longrightarrow\mathcal P^n_{X/S}.\tag{4.6}
\]
根据\Term{isomorphism}{同构} (4.4)，这个\Term{homomorphism}{同态}把 $\mathcal O_X$ 的\Term{section}{截面} $f$ 映成 $\mathcal P^n_{X/S}$ 的\Term{section}{截面} $d^n(f)$；后者沿 $X$ 在 $S$ 上的任一\Term{section}{截面} $s$ 的\Term{inverse-image}{逆像}，等同于 $f$ 在 $\Gamma(X,\mathcal O_X/\mathcal J^{n+1})$ 中的典范像。这说明，把 $d^n f$ 称为“$f$ 的 $n$ 阶\Term{principal-parts}{主部系}”是恰当的；公式 (4.5) 中 $S=\operatorname{Spec}(k)$ 的情形尤为如此。

最后注意，\Term{homomorphism}{同态} (4.6) 可以看成 $\mathcal O_X$ 上相对于常数\Term{prescheme}{预概形} $S$ 的普遍的阶数不超过 $n$ 的\Term{differential-operator}{微分算子}\footnote{本段所涉及的内容可参阅 EGA IV 16.8 至 16.12。}。这里约定，从 $\mathcal O_X$ 到模 $F$ 的阶数不超过 $n$ 的\Term{differential-operator}{微分算子}，是指\Term{separable}{可分}解为
\[
D:\mathcal O_X\xrightarrow{\ d^n\ }\mathcal P^n_{X/S}\xrightarrow{\ u\ }F
\]
的\Term{sheaf}{层}\Term{homomorphism}{同态} $D$，其中 $u$ 是 $\mathcal O_X$-模\Term{homomorphism}{同态}，而且由 $D$ 唯一确定。这个定义与直观的递归定义一致：$D$ 是阶数不超过 $n$ 的\Term{differential-operator}{微分算子}，是指对 $X$ 的任意开集 $U$ 上的 $\mathcal O_X$-\Term{section}{截面} $g$，映射 $f\mapsto D(fg)-D(f)$ 是 $U$ 上阶数不超过 $n-1$ 的\Term{differential-operator}{微分算子}。因此，若 $X$ 在 $S$ 上\Term{differentially-smooth}{微分光滑}，则所有阶数的\Term{differential-operator}{微分算子}组成的环\Term{sheaf}{层}具有可微流形上的微分计算中熟知的简单结构。特别地，它局部有一组 $\mathcal O_X$-模的基，由两两可交换的算子 $\delta/\delta x_i$（$1\leq i\leq n$）的\Term{divided-power}{除幂}构成。

\SourcePage{59}
若 $S$ 是一个 $\mathbb Q$-代数\Term{sheaf}{层}，其中 $\mathbb Q$ 为有理数域，则只须取各 $\delta/\delta x_i$ 的普通多项式。在这个非常特殊的情形，要使 $X$ 在 $S$ 上\Term{differentially-smooth}{微分光滑}，只要 $\Omega^1_{X/S}$ \Term{locally-free}{局部自由}就够了。

\textbf{注记 4.19。}本节引入的“\Term{regular-immersion}{正则浸入}”“\Term{regular-ideal}{正则理想}”等术语，曾受到颇为强烈而广泛的反对，Chevalley、Serre 都在其列。有人提议说“Cohen--Macaulay 理想”“Macaulay 理想”或“麦考利型理想”；照此便也应当采用“Cohen--Macaulay \Term{immersion}{浸入}”或“Macaulay \Term{immersion}{浸入}”。然而，这与巨著未来稿件中已经使用的另一术语冲突：那里，一个\Term{finite-type}{有限型}\Term{morphism}{态射}称为在一点处“\Term{cohen-macaulay}{Cohen--Macaulay}”，是指它在该点\Term{flat}{平坦}，并且经过该点的\Term{fiber}{纤维}在该点的\Term{local-ring}{局部环}是\Term{cohen-macaulay-ring}{Cohen--Macaulay 环}。在找到满意的解决办法以前，我们暂且保留本节引入的术语，但仍持保留态度。\footnote{这正是 EGA $0_{\mathrm{IV}}$ 15.1.7 采用的术语。}

\SGAsection{5}{基底为域的情形}
\begin{SGAresult}{命题}{5.1}
设 $k$ 是域，$X$ 是 $k$ 上\Term{finite-type}{有限型}的\Term{prescheme}{预概形}，$x\in X$，$n$ 是 $X$ 在 $x$ 处的维数；设
\[
f:X\longrightarrow\operatorname{Spec}k[t_1,\ldots,t_n]=Y
\]
是由元素 $f_i\in\Gamma(X,\mathcal O_X)$ 定义的\Term{morphism}{态射}。下列条件等价，而且都蕴含 $X$ 在 $x$ 处在 $k$ 上\Term{smooth}{光滑}，从而由 3.1，$X$ 在 $x$ 处\Term{regular}{正则}：
\begin{enumerate}
\item[(i)] $f$ 在 $x$ 处\Term{etale}{平展}；
\item[(ii)] 各 $df_i$ 在 $x$ 处构成 $\Omega^1_{X/k}$ 的一组基；
\item[(iii)] 各 $df_i$ 在 $x$ 处生成 $\Omega^1_{X/k}$。
\end{enumerate}
\end{SGAresult}
因为 (i) 蕴含 $X$ 在 $x$ 处在 $k$ 上\Term{smooth}{光滑}，所以 (i)$\Rightarrow$(ii) 是 4.8 的一个特例。(ii)$\Rightarrow$(iii) 显然，只须证明 (iii)$\Rightarrow$(i)。若 (iii) 成立，由引理 4.1，$f$ 在 $x$ 处\Term{net}{净}。于是以 $x$ 的某个\Term{open-neighborhood}{开邻域}替换 $X$ 后，$f$ \Term{quasi-finite}{拟有限}，再由维数可知它\Term{dominant}{支配}。由于 $Y$ \Term{regular}{正则}，由 I 9.5 (ii) 或 I 9.11，$f$ \Term{etale}{平展}。

\begin{SGAresult}{推论}{5.2}
在 5.1 的预备条件下，再设 $k(x)$ 是 $k$ 的\Term{finite}{有限}\Term{separable-extension}{可分扩张}，且 $f_i$（$1\leq i\leq n$）在 $\mathfrak m_x$ 中定义元素。则上述条件还等价于：
\begin{enumerate}
\item[(iv)] 各 $f_i$ 构成 $\mathfrak m_x$ 的一组生成元；等价地，各 $f_i\bmod\mathfrak m_x^2$ 构成 $\mathfrak m_x/\mathfrak m_x^2$ 的一组 $k(x)$-基。
\end{enumerate}
\end{SGAresult}
确实，利用\Term{exact-sequence}{正合列}
\[
\mathfrak m_x/\mathfrak m_x^2\longrightarrow\Omega^1_{\mathcal O_x/k}
\longrightarrow\Omega^1_{k(x)/k}\longrightarrow0,\tag{5.1}
\]
\SourcePage{60}
并注意到 $k(x)$ 在 $k$ 上\Term{etale}{平展}，所以 $\Omega^1_{k(x)/k}=0$，即知 (iv)$\Rightarrow$(iii)。另一方面，(ii) 蕴含 (iv)，因为 $X$ 与 $\operatorname{Spec}(k(x))$ 都在 $x$ 处在 $k$ 上\Term{smooth}{光滑}，故由 4.10 (iv)，可在上述\Term{exact-sequence}{正合列}左端补上一个 $0$。

\begin{SGAresult}{推论}{5.3}
设 $x$ 是 $k$ 上\Term{finite-type}{有限型}的 $X$ 的一点。若 $X$ 在 $x$ 处在 $k$ 上\Term{smooth}{光滑}，则 $\mathcal O_x$ \Term{regular}{正则}。若 $k(x)$ 是 $k$ 的\Term{finite}{有限}\Term{separable-extension}{可分扩张}，则逆命题成立。
\end{SGAresult}
逆命题确实由 5.2 得出，只须取 $\mathfrak m_x$ 的一组\Term{regular-generating-system}{正则生成元系} $(f_i)$。注意，也可以用定理 4.15 代替 5.2。由此得到：

\begin{SGAresult}{命题}{5.4}
设 $X$ 是 $k$ 上\Term{finite-type}{有限型}的\Term{prescheme}{预概形}。若 $X$ 在 $k$ 上\Term{smooth}{光滑}，则 $X$ \Term{regular}{正则}。若 $k$ 是\Term{perfect-field}{完美域}，则逆命题成立。
\end{SGAresult}
对于逆命题，由 5.3，$X$ 在每个\Term{closed-point}{闭点}都在 $k$ 上\Term{smooth}{光滑}；又因\Term{smooth}{光滑}点的集合是开集，故处处\Term{smooth}{光滑}。

\begin{SGAresult}{定理}{5.5}
设 $X$ 是 $k$ 上\Term{finite-type}{有限型}的\Term{prescheme}{预概形}，$x\in X$，$n$ 是 $X$ 在 $x$ 处的维数，$k'$ 是 $k$ 的一个\Term{perfect}{完美}\Term{group-extension}{扩张}。下列条件等价：
\begin{enumerate}
\item[(i)] $X$ 在 $x$ 处在 $k$ 上\Term{smooth}{光滑}；
\item[(ii)] $\Omega^1_{X/k}$ 在 $x$ 处自由，\Term{rank}{秩}为 $n$；
\item[(ii bis)] $\Omega^1_{X/k}$ 在 $x$ 处可由 $n$ 个元素生成；
\item[(iii)] $X$ 在 $x$ 处在 $k$ 上\Term{differentially-smooth}{微分光滑}；
\item[(iv)] $x$ 有一个\Term{open-neighborhood}{开邻域} $U$，使 $U\otimes_k k'$ \Term{regular}{正则}，即其各点的\Term{local-ring}{局部环}均\Term{regular}{正则}。
\end{enumerate}
\end{SGAresult}
由 4.3 (ii)，有 (i)$\Rightarrow$(ii)；(ii)$\Rightarrow$(ii bis) 显然；(ii bis)$\Rightarrow$(i) 由 5.1 得出。由于 $X$ 在 $k$ 上\Term{flat}{平坦}，由 4.18，(i)$\Leftrightarrow$(iii)。\Term{smoothness}{光滑性}在\Term{base-change}{基变换}下保持，且蕴含\Term{regularity}{正则性}，故 (i)$\Rightarrow$(iv)。反过来，由命题 5.4，$U\otimes_k k'$ 在 $k'$ 上简单，再由 4.13，$U$ 在 $k$ 上简单，故 (iv)$\Rightarrow$(i)。

设 $X$ \Term{irreducible}{不可约}，取 $x$ 为它的\Term{generic-point}{泛点}，得到：
\begin{SGAresult}{推论}{5.6}
设 $K$ 是一个局部\Term{artinian-ring}{阿廷环}，它是域 $k$ 上某个\Term{finite-type}{有限型}代数的\Term{localization}{局部化}；例如，$K$ 可以是 $k$ 的\Term{finitely-generated-extension}{有限生成扩张}。令 $n$ 为其\Term{residue-field}{剩余域}在 $K$ 上的\Term{transcendence-degree}{超越次数}。下列条件等价：
\begin{enumerate}
\item[(i)] $K$ 是 $k$ 的某个\Term{purely-transcendental-extension}{纯超越扩张} $k(t_1,\ldots,t_n)$ 的\Term{finite}{有限}\Term{separable-extension}{可分扩张}；
\item[(ii)] $\Omega^1_{X/k}$ 是\Term{rank}{秩}为 $n$ 的自由 $K$-模；
\item[(ii bis)] $\Omega^1_{X/k}$ 是可由 $n$ 个元素生成的 $K$-模；
\item[(iii)] $K\otimes_k K$ 关于增广 $K\otimes_k K\to K$ 的核的幂所定义拓扑的\Term{completion}{完备化} $\mathcal O'$，是一个“\Term{regular}{正则}”的\Term{augmented-algebra}{增广} $K$-代数，即\Term{isomorphism}{同构}于一个以 $K$ 为系数的\Term{formal-power-series-algebra}{形式幂级数代数}。注意，若 $K$ 是域，这等价于说 $\mathcal O'$ 是\Term{regular-local-ring}{正则局部环}；
\SourcePage{61}
\item[(iv)] $K$ 是 $k$ 的\Term{separable-extension}{可分扩张}。
\end{enumerate}
\end{SGAresult}
确实，总可把 $K$ 看成 $k$ 上某个\Term{irreducible}{不可约}\Term{finite-type}{有限型}\Term{scheme}{概形} $X$ 的\Term{generic-point}{泛点}的\Term{local-ring}{局部环}。所考虑的各条件就是 5.5 中同名的条件，其中 (iv) 取 $k'$ 为 $k$ 的一个\Term{algebraically-closed}{代数闭}\Term{group-extension}{扩张}。只须证明“$K$ 在 $k$ 上\Term{separable}{可分}”蕴含“$X$ 在 $x$ 处在 $k$ 上\Term{smooth}{光滑}”。利用 4.13，可立即归结为基域为 $k'$ 的情形，于是基域\Term{algebraically-closed}{代数闭}，因而 $X$ 有一个 $k$-有理点 $a$。由 5.4，$X$ 在 $a$ 处在 $k$ 上\Term{smooth}{光滑}，从而在\Term{generic-point}{泛点} $x$ 处也\Term{smooth}{光滑}，证毕。\footnote{见本讲末尾的勘误，原书重排本第 47 页。}

注意，当 $K$ 是 $k$ 的\Term{finitely-generated-extension}{有限生成扩张}时，(i)、(ii)、(ii bis)、(iv) 的等价性早已熟知，但这里没有使用这些已知等价性中的任何一个。当然，命题 5.1 还包含如下特例：$K$ 中一组元素 $x_i$（$1\leq i\leq n$）是 $K$ 在 $k$ 上的一组“\Term{separating-transcendence-basis}{可分超越基}”，当且仅当各 $dx_i$ 构成 $K$-模 $\Omega^1_{K/k}$ 的一组基。这也是一个熟知的事实。

\begin{SGAresult}{推论}{5.7}
设 $X$ 是域 $k$ 上\Term{finite-type}{有限型}的\Term{prescheme}{预概形}。$X$ 在 $k$ 上\Term{smooth}{光滑}的充要条件是：$\Omega^1_{X/k}$ \Term{locally-free}{局部自由}，且 $X$ 的各\Term{irreducible-component}{不可约分支}的\Term{generic-point}{泛点}的\Term{local-ring}{局部环}都是 $k$ 的\Term{separable-extension}{可分扩张}。若 $k$ \Term{perfect}{完美}且 $X$ \Term{reduced}{约化}，则最后这个条件自动成立。
\end{SGAresult}
可设 $X$ \Term{connected}{连通}，令 $n$ 为已假设\Term{locally-free}{局部自由}的 $\Omega^1_{X/k}$ 的\Term{rank}{秩}。由假设和 5.6，$n$ 也是由 $X$ 的各\Term{generic-point}{泛点}的\Term{local-ring}{局部环}定义的 $k$-\Term{group-extension}{扩张}的\Term{transcendence-degree}{超越次数}，所以 $X$ 的所有\Term{irreducible-component}{不可约分支}都为 $n$ 维。再由 5.5 得出结论。

须注意，若 $K$ 是 $k$ 的\Term{finite}{有限}\Term{group-extension}{扩张}，不必\Term{separable}{可分}，则 $\Omega^1_{K/k}$ 是自由 $k$-模。故令 $X=\operatorname{Spec}(K)$，则 $\Omega^1_{X/k}$ 是\Term{locally-free}{局部自由}\Term{sheaf}{层}，$X$ \Term{reduced}{约化}，但 $X$ 未必在 $k$ 上\Term{smooth}{光滑}。再把标量\Term{group-extension}{扩张}到 $k$ 的\Term{algebraic-closure}{代数闭包}，便得到一个类似例子：此时 $k$ \Term{algebraically-closed}{代数闭}，但 $X$ 不再\Term{reduced}{约化}。

\begin{SGAresult}{推论}{5.8}
设 $X$ 是域 $k$ 上\Term{finite-type}{有限型}的\Term{prescheme}{预概形}，$x\in X$，$n$ 是 $X$ 在 $x$ 处的维数，$p$ 是 $\mathcal O_x$ 的维数，即 $x$ 在 $X$ 中的闭包 $Y$ 的\Term{codimension}{余维数}。于是 $n-p$ 是 $k(x)$ 在 $k$ 上的\Term{transcendence-degree}{超越次数}。设 $f_i$（$1\leq i\leq n$）是 $\mathcal O_x$ 中的元素，且当 $1\leq i\leq p$ 时有 $f_i\in\mathfrak m_x$。下列条件等价：
\begin{enumerate}
\item[(i)] 由各 $f_i$ 定义的在 $x$ 处的\Term{morphism}{态射}芽
\[
X\longrightarrow\operatorname{Spec}(k[t_1,\ldots,t_n])
\]
在 $x$ 处\Term{etale}{平展}。
\SourcePage{62}
\item[(ii)] 各 $f_i$（$1\leq i\leq p$）生成 $\mathfrak m_x$，即构成 $\mathcal O_x$ 的一组\Term{regular-system-of-parameters}{正则参数系}；各 $f_j$（$p+1\leq j\leq n$）在 $k(x)$ 中的类构成\Term{separating-transcendence-basis}{可分超越基}，即各 $d\overline f_j$（$p+1\leq j\leq n$）构成 $\Omega^1_{k(x)/k}$ 的一组基，或等价地，生成 $\Omega^1_{k(x)/k}$。
\end{enumerate}
\end{SGAresult}
\begin{proof}
设 (i) 成立。由 4.8，各 $df_i(x)$ 构成 $\Omega^1_{X/k}(x)$ 的一组基，故它们在 $\Omega^1_{k(x)/k}$ 中的像 $d\overline f_i(x)$ 在 $k$ 上生成这个向量空间。因为 $\overline f_i=0$（$1\leq i\leq p$），所以只须取 $p+1\leq i\leq n$ 的 $d\overline f_i(x)$。由于 $k(x)$ 在 $k$ 上的\Term{transcendence-degree}{超越次数}为 $n-p$，由推论 5.6 的准则 (iii)，应用于 $K=k(x)$，可知 $Y$ 在其\Term{generic-point}{泛点} $x$ 处在 $k$ 上\Term{smooth}{光滑}，且各 $d\overline f_i(x)$（$p+1\leq i\leq n$）构成 $\Omega^1_{k(x)/k}$ 的一组 $k(x)$-基。因此满足 4.9 的条件 (ii)，从而也满足条件 (iii)。特别地，各 $f_i$（$1\leq i\leq p$）构成 $\mathfrak m_x$ 的一组生成元。由于 $\mathcal O_x$ 的维数是 $p$，它们构成 $x$ 处的一组\Term{regular-system-of-parameters}{正则参数系}。这就证明了 (ii)。

设 (ii) 成立。由\Term{exact-sequence}{正合列} (5.1)，各 $df_i(x)$ 生成 $\Omega^1_{X/k}$，再由命题 5.1 得到 (i)。
\end{proof}

\begin{SGAresult}{推论}{5.9}
设 $X$ 是域 $k$ 上\Term{finite-type}{有限型}的\Term{prescheme}{预概形}，$x\in X$，$n$ 是 $X$ 在 $x$ 处的维数，$p$ 是 $\mathcal O_x$ 的维数，即 $x$ 在 $X$ 中的闭包 $Y$ 的\Term{codimension}{余维数}；于是 $n-p$ 为 $k(x)$ 在 $k$ 上的\Term{transcendence-degree}{超越次数}。下列条件等价：
\begin{enumerate}
\item[(i)] $\mathcal O_x$ \Term{regular}{正则}，且 $k(x)$ 是 $k$ 的\Term{separable-extension}{可分扩张}。
\item[(ii)] $X$ 在 $x$ 处在 $k$ 上\Term{smooth}{光滑}，且典范\Term{homomorphism}{同态}
\[
\mathfrak m_x/\mathfrak m_x^2\longrightarrow
\Omega^1_{\mathcal O_x/k}\otimes_{\mathcal O_x}k(x)=\Omega^1_{X/k}(x)
\]
为\Term{injectivity}{单射}。
\item[(iii)] 存在 $f_i\in\mathcal O_x$（$1\leq i\leq n$），当 $1\leq i\leq p$ 时 $f_i\in\mathfrak m_x$，使它们定义的从 $X$ 到 $\operatorname{Spec}(k[t_1,\ldots,t_n])$ 的\Term{morphism}{态射}芽在 $x$ 处\Term{etale}{平展}；即由 5.1，各 $df_i(x)$ 生成 $\Omega^1_{X/k}(x)$。
\item[(iv)] 存在 $f_i\in\mathcal O_x$（$1\leq i\leq n$），使 $f_i$（$1\leq i\leq p$）生成 $\mathfrak m_x$，且各 $d\overline f_j(x)$（$p+1\leq j\leq n$）在 $k(x)$ 上生成 $\Omega^1_{k(x)/k}$。
\end{enumerate}
\end{SGAresult}
(iii)、(iv) 的等价性来自推论 5.8；由 4.9，这些条件还等价于 $X$ 在 $x$ 处在 $k$ 上\Term{smooth}{光滑}，且 4.10 的条件 (ii) 成立。因此它们等价于 $X$ 在 $x$ 处在 $k$ 上\Term{smooth}{光滑}，且 4.10 的条件 (iv) 成立，也就等价于 5.9 (ii)。它们也等价于 $X$ 在 $x$ 处在 $k$ 上\Term{smooth}{光滑}，且 4.10 的条件 (i) 成立；在这里，后一个条件只是说 $k(x)$ 在 $k$ 上\Term{separable}{可分}。这蕴含 5.9 (i)。只须再证明 5.9 (i) 蕴含上述条件，也就是证明下一个推论。

\SourcePage{63}
\begin{SGAresult}{推论}{5.10}
设 $x$ 是域 $k$ 上\Term{finite-type}{有限型}的\Term{prescheme}{预概形} $X$ 的一点，且 $k(x)$ 在 $k$ 上\Term{separable}{可分}。则 $X$ 在 $x$ 处在 $k$ 上\Term{smooth}{光滑}的充要条件是：$X$ 在 $x$ 处\Term{regular}{正则}，即\Term{local-ring}{局部环} $\mathcal O_x$ \Term{regular}{正则}。
\end{SGAresult}
确实，若满足\Term{regularity}{正则性}，显然可取满足条件 5.9 (iv) 的 $f_i\in\mathcal O_x$（$1\leq i\leq n$）。

\textbf{勘误。}本节 5.6 的证明使用了如下事实：\Term{algebraically-closed-field}{代数闭域}上的非空\Term{reduced}{约化}\Term{finite-type}{有限型}\Term{scheme}{概形}至少有一个\Term{regular-point}{正则点}，因而有一个\Term{smooth-point}{光滑点}。通常借助 Zariski 关于 $X$ 的\Term{regular-point}{正则点}集为开集的定理，以微分方法证明这一事实。为避免循环论证，须证明：若 $K/k$ 是\Term{finitely-generated-extension}{有限生成}的\Term{separable-extension}{可分扩张}，且 $f_i\in K$ 满足各 $d_{K/k}f_i$（$1\leq i\leq n$）构成 $\Omega^1_{K/k}$ 的一组基，则 $n$ 就是 $K$ 在 $k$ 上的\Term{transcendence-degree}{超越次数}，即各 $f_i$ \Term{algebraically-independent}{代数无关}。用 Mac Lane 准则证明这个事实的方法是熟知的，见 Bourbaki《代数》第 V 章第 9 节定理 2。取次数最小且满足 $g(f_1,\ldots,f_n)=0$ 的多项式 $g\in k[t_1,\ldots,t_n]$，则
\[
\sum_i\frac{dg}{dt_i}(f_1,\ldots,f_n)\,df_i=0.
\]
由于各 $df_i$ 构成 $\Omega^1_{K/k}$ 的一组基，可知各 $dg/dt_i$ 在 $(f_1,\ldots,f_n)$ 处取值为零。由 $g$ 的次数极小性，它们都是零多项式。若 $k$ 的特征为 $0$，则 $g=0$；若 $k$ 的特征为 $p\ne0$，则 $g=h(t_1^p,\ldots,t_n^p)$。利用 Mac Lane 准则，可见多项式 $h\in k[t_1,\ldots,t_n]$ 也在 $(f_1,\ldots,f_n)$ 处取值为零，再由 $g$ 的次数极小性，仍得 $g=0$。
% PDF 页 64 为空白，没有待译正文。
